% =============================================================================
% BOOK: Cooperative Magnetic ASV Search — System & ROS 2 Architecture Manual
% Overleaf: set compiler to pdfLaTeX; main file = main.tex
% Upload the entire overleaf_magnetic_asv/ folder (including chapters/).
% =============================================================================
\documentclass[11pt,a4paper,oneside,openany]{book}

\usepackage[margin=2.4cm]{geometry}
\usepackage{amsmath,amssymb,amsthm}
\usepackage{booktabs}
\usepackage{graphicx}
\usepackage{hyperref}
\usepackage{xcolor}
\usepackage{setspace}
\usepackage{enumitem}
\usepackage{caption}
\usepackage{longtable}
\usepackage{array}
\usepackage{multirow}
\usepackage{float}
\usepackage{fancyhdr}
\usepackage{csquotes}
\usepackage[numbers,sort&compress]{natbib}

\hypersetup{
  colorlinks=true,
  linkcolor=blue!45!black,
  citecolor=blue!45!black,
  urlcolor=blue!55!black
}

\pagestyle{fancy}
\fancyhf{}
\fancyhead[L]{\small\nouppercase{\leftmark}}
\fancyhead[R]{\small Magnetic ASV --- ROS~2 Manual}
\fancyfoot[C]{\thepage}
\renewcommand{\headrulewidth}{0.4pt}

\title{\textbf{Cooperative Magnetic Anomaly Search\\
with Differential-Thrust ASVs}\\[0.6em]
\large A Book-Style Technical Manual:\\
Algorithms, Pipelines, and Complete ROS~2 Architecture}
\author{
  System Documentation --- Competition / Proposal Package\\
  ROS~2 Jazzy $\cdot$ Gazebo Harmonic $\cdot$ \texttt{version\_3} three-ASV strip stack
}
\date{\today}

\begin{document}
\frontmatter
\maketitle

\chapter*{Preface}
\addcontentsline{toc}{chapter}{Preface}
This manual is written like a software architecture book, not a short paper
abstract. The goal is that a judge, teammate, or new developer can open a
chapter, pick a ROS node, and know exactly:

\begin{itemize}[leftmargin=*]
  \item which package owns it and how it is launched;
  \item every parameter (code default vs YAML override);
  \item every subscription and publication (relative vs absolute topic);
  \item the callback / timer logic in order;
  \item how that node sits in the end-to-end pipeline.
\end{itemize}

The reference implementation described here is the three-ASV
\texttt{version\_3} workspace (\(300\,\mathrm{m}\) lake, strip coverage). Larger
lakes reuse the same nodes with scaled configuration. See also the Markdown
developer manual \texttt{docs/MAGNETIC\_ASV\_ROS2\_DEVELOPER\_MANUAL.md}.

\tableofcontents

\mainmatter

% ---------------------------------------------------------------------------
\chapter{Introduction and document map}
\label{ch:intro}

\section{What this system does}
We search a bounded lake for a compact magnetic dipole anomaly using one or
more differential-thrust autonomous surface vehicles (ASVs). Each boat runs an
identical ROS~2 node graph under its own namespace (\texttt{/asv1},
\texttt{/asv2}, \texttt{/asv3}). A shared \texttt{/swarm} namespace hosts
Bayesian fusion and verification arbitration. Gazebo Harmonic provides
hydrodynamics and sensor plugins; \texttt{ros\_gz\_bridge} and a
Gazebo-Transport pose helper connect the plant to ROS.

The reference stack for this book is \texttt{version\_3}: a
\(300\times300\,\mathrm{m}\) lake with up to \textbf{three ASVs}, vertical
\textbf{strip} lawnmower coverage (dead corridor between strips), and peer
verification after spiral densification. A long-form Markdown companion for
developers lives at
\texttt{docs/MAGNETIC\_ASV\_ROS2\_DEVELOPER\_MANUAL.md}.

\section{How to read this book}
\begin{description}[style=nextline,leftmargin=1.2cm]
  \item[Chapter~\ref{ch:pipelines}] End-to-end pipelines: control loop, sensing
        chain, belief loop, mission loop. Read this first to see how nodes
        chain.
  \item[Chapter~\ref{ch:algorithms}] Mathematical models (dipole, Bayes, MI,
        inland-safe spiral, LOS).
  \item[Chapter~\ref{ch:fsm}] Mission finite-state machine and cooperative
        verify handshake.
  \item[Chapter~\ref{ch:ros-catalog}] \textbf{The core of this manual:} every
        message, bridge, and node with parameters, topics, and logic.
  \item[Chapter~\ref{ch:ops}] Launches, namespaces, bags, and operator UI.
  \item[Chapter~\ref{ch:accept}] What ``confirmed'' means in software vs
        offline scoring against the planted target.
\end{description}

\section{Package inventory}
\begin{center}
\begin{tabular}{@{}lp{9.5cm}@{}}
\toprule
Package & Role \\
\midrule
\texttt{boat\_bringup} & Launches, YAML, Gazebo bridges \\
\texttt{boat\_description} & SDF boats + lake world \\
\texttt{boat\_msgs} & Custom message contracts \\
\texttt{boat\_control} & \texttt{thrust\_mixer} \\
\texttt{boat\_navigation} & Pose, LOS follower, plotter \\
\texttt{boat\_sensing} & Mag driver, filter, calibration \\
\texttt{boat\_mapping} & Bayes fusion + dipole fit \\
\texttt{boat\_mission} & Mission FSM, planning, verify coordinator \\
\bottomrule
\end{tabular}
\end{center}

\section{Namespace rules (critical)}
\begin{itemize}[leftmargin=*]
  \item \textbf{Relative topics} (no leading slash), e.g.\ \texttt{cmd\_vel},
        resolve under the node namespace: on \texttt{/asv1} they become
        \texttt{/asv1/cmd\_vel}.
  \item \textbf{Absolute topics} (leading slash), e.g.\
        \texttt{/swarm/belief/peak} or
        \texttt{/model/simple\_boat/.../cmd\_thrust}, are global and identical
        for every ASV.
  \item Per-ASV stacks are launched \emph{inside} \texttt{/asv\(k\)}. Swarm
        nodes (\texttt{bayes\_fusion}, \texttt{verify\_coordinator}) are
        launched at the root.
  \item Model map: \texttt{asv1}$\leftrightarrow$\texttt{simple\_boat},
        \texttt{asv2}$\leftrightarrow$\texttt{simple\_boat\_2},
        \texttt{asv3}$\leftrightarrow$\texttt{simple\_boat\_3}.
\end{itemize}

\section{Primary demo}
\begin{verbatim}
cd version_3 && source install/setup.bash
ros2 launch boat_bringup multi_asv_phase4.launch.py
\end{verbatim}
Always kill leftover Gazebo/ROS/\texttt{thrust\_mixer} processes before a new
run; latched \texttt{/swarm/mission/complete} from a prior domain looks like a
false CONFIRMED.

\chapter{End-to-end ROS pipelines}
\label{ch:pipelines}

This chapter walks the live system as pipelines. Each arrow is a ROS topic
(or Gazebo Transport channel). Chapter~\ref{ch:ros-catalog} then expands every
node.

\section{Pipeline A --- Motion control (closed loop)}
\begin{verbatim}
Gazebo dynamics
   |  (Gazebo Transport: /world/.../dynamic_pose/info)
   v
gazebo_pose2d  --->  /asvK/pose2d
   |
   |  also: ros_gz_bridge ---> /asvK/odom
   v
los_path_follower
   subs: /asvK/odom , /asvK/plan   (plan is TRANSIENT_LOCAL)
   pubs: /asvK/cmd_vel , /asvK/plan/active , /asvK/nav/debug
   v
thrust_mixer
   subs: /asvK/cmd_vel
   pubs: /model/<boat>/joint/left|right_propeller_joint/cmd_thrust
   v
ros_gz_bridge (ROS_TO_GZ) ---> Gazebo propeller joints
\end{verbatim}

\paragraph{Who owns the path?}
\texttt{mission\_manager} (or \texttt{spiral\_demo}) publishes
\texttt{/asv\(K\)/plan}. The LOS follower latches the latest Path (QoS
transient local) and tracks it until a new plan arrives. During HOLD / COMPLETE
the mission manager publishes a trivial hold path near the current pose.

\paragraph{Logic order inside one control cycle.}
\begin{enumerate}[leftmargin=*]
  \item Odometry updates boat pose for LOS geometry.
  \item Cross-track error and lookahead yield a commanded heading.
  \item Yaw PID $\rightarrow$ \texttt{angular.z}; surge policy $\rightarrow$
        \texttt{linear.x}.
  \item Mixer maps \((v,\omega)\) to left/right thrust floats.
  \item Bridge writes thrusts into Gazebo; hydrodynamics integrate.
\end{enumerate}

\section{Pipeline B --- Magnetometry to cleaned anomaly}
\begin{verbatim}
Gazebo magnetometer
   v
ros_gz_bridge ---> /asvK/mag/gazebo   (sensor_msgs/MagneticField, Tesla)
   v
mag_driver
   also uses: /asvK/pose2d , /asvK/cmd_vel
   optionally overlays planted dipole A/(r^3+s^3)
   pubs: /asvK/mag/raw                 (boat_msgs/MagReading, nT)
   v
mag_filter
   spike reject + moving-average LPF
   pubs: /asvK/mag/filtered , /asvK/mag/filter_status
   v
calibration_node
   spatial heading-binned baseline + temporal high-pass
   pubs: /asvK/mag/anomaly             (boat_msgs/MagAnomaly)
         /asvK/calibration/status      (CALIBRATING|READY|DEGRADED)
\end{verbatim}

\paragraph{Why three stages?}
\begin{itemize}[leftmargin=*]
  \item \texttt{mag\_driver} owns units, pose stamping, motor flag, and the
        synthetic dipole plant used only in simulation evaluation.
  \item \texttt{mag\_filter} owns robust denoising so calibration does not
        learn spikes.
  \item \texttt{calibration\_node} owns ambient baseline learning and the
        cleaned anomaly feature consumed by Bayes and verification.
\end{itemize}

\section{Pipeline C --- Shared belief and localization}
\begin{verbatim}
/asv1/mag/anomaly  --+
                     +-->  bayes_fusion  -->  /swarm/belief/map
/asv2/mag/anomaly  --+                    -->  /swarm/belief/peak
                                          -->  /swarm/belief/peak_probability
                                          -->  /swarm/belief/centroid
                                          -->  /swarm/belief/centroid_mass
                                          -->  /swarm/belief/centroid_spread
                                          -->  /swarm/belief/fix
                                          -->  /swarm/belief/fix_rms
                                          -->  /swarm/belief/fix_samples
\end{verbatim}

\paragraph{Internal order in \texttt{bayes\_fusion}.}
\begin{enumerate}[leftmargin=*]
  \item Receive \texttt{MagAnomaly}; ignore if not calibrated (ABSTAIN).
  \item Classify HIT / MISS / ABSTAIN against thresholds.
  \item Update every grid cell with \(P_{\mathrm{det}}(d)\) likelihood; renormalize.
  \item Compute discrete peak and belief-weighted centroid (+ spread).
  \item Buffer strong samples into dipole fitter; optionally publish continuous fix.
\end{enumerate}

\section{Pipeline D --- Mission behaviour and swarm verify}
\begin{verbatim}
mission_manager (/asvK)
   subs: pose2d, mag/anomaly, calibration/status,
         /swarm/belief/{map,peak,peak_probability,centroid},
         /swarm/verify/request , /swarm/mission/halt
   pubs: /asvK/plan , /asvK/mission/state , /asvK/mission/mode ,
         /asvK/mission/info_gain ,
         /swarm/verify/declare , /swarm/verify/result

verify_coordinator (global, once)
   subs: /swarm/verify/declare , /swarm/verify/result
   pubs: /swarm/verify/request , /swarm/mission/complete ,
         /swarm/mission/halt , /swarm/mission/status
\end{verbatim}

\paragraph{Dual-ASV happy path.}
\begin{enumerate}[leftmargin=*]
  \item Both boats lawnmower Voronoi regions (\texttt{GLOBAL\_SEARCH}).
  \item Belief HITs raise \(p^\star\); nearby boat enters
        \texttt{TARGET\_SEARCH} (MI then spiral).
  \item Discoverer declares candidate on \texttt{/swarm/verify/declare} and
        HOLDs.
  \item Coordinator assigns the other ASV on \texttt{/swarm/verify/request}.
  \item Verifier orbits, counts strong confirmations, publishes
        \texttt{VerifyResult}.
  \item Coordinator latches COMPLETE / HALT; plotter banner shows
        \textbf{CONFIRMED}.
\end{enumerate}

\section{Pipeline E --- Operator visualization}
\texttt{trajectory\_plotter} (global in multi-ASV mode) subscribes to every
boat's \texttt{pose2d}, \texttt{plan/active}, \texttt{mag/anomaly},
\texttt{mission/state}, plus all \texttt{/swarm/belief/*} and verify/complete
topics. It does not publish commands. The bottom status strip mirrors mission
phase chips: SEARCH $\rightarrow$ HUNT $\rightarrow$ VERIFY $\rightarrow$
CONFIRMED.

\section{Minimum viable process set}
To move one boat under LOS alone you need: Gazebo, bridges,
\texttt{gazebo\_pose2d}, \texttt{los\_path\_follower}, \texttt{thrust\_mixer},
and someone publishing a Path. To run the full magnetic mission you additionally
need the sensing trio, \texttt{bayes\_fusion}, \texttt{mission\_manager}, and
(for cooperative verify) \texttt{verify\_coordinator}.

\chapter{Mathematical algorithms}
\label{ch:algorithms}

\section{Planted dipole (simulation evaluation only)}
\[
  a(\mathbf{p})=\frac{A}{r^{3}+s^{3}},
  \qquad r=\lVert\mathbf{p}-\mathbf{t}\rVert_{2}.
\]
Defaults in \texttt{sensing.yaml}: \(A=4\times10^{5}\,\mathrm{nT\cdot m^{3}}\),
\(s=20\,\mathrm{m}\) $\Rightarrow$ \(50\,\mathrm{nT}\) overhead. Current planted
pose in \texttt{version\_3}: \(\mathbf{t}=(85,40)\,\mathrm{m}\). Estimators
never read \(\mathbf{t}\); they see only cleaned anomalies.

\section{Detection probability}
\[
  P_{\mathrm{det}}(d)=P_{\mathrm{bg}}+(P_{\max}-P_{\mathrm{bg}})
  \frac{d_{1/2}^{3}}{d^{3}+d_{1/2}^{3}}.
\]
Defaults: \(P_{\mathrm{bg}}=0.05\), \(P_{\max}=0.95\), \(d_{1/2}=30\,\mathrm{m}\).

\section{HIT / MISS / ABSTAIN}
With cleaned anomaly \(a_{\mathrm{cl}}\): HIT if
\(a_{\mathrm{cl}}\ge\tau_{\mathrm{hit}}\) (default \(15\,\mathrm{nT}\)); MISS if
\(\le\tau_{\mathrm{miss}}\) (default \(5\,\mathrm{nT}\)); else ABSTAIN.
\texttt{hit\_only: true} treats MISS as ABSTAIN.

\section{Bayesian grid update}
Uniform prior on an \(N\times N\) grid (\(\Delta=10\,\mathrm{m}\),
\(300\,\mathrm{m}\) map, origin \((-150,-150)\)). On HIT,
\(b_k\leftarrow b_k P_{\mathrm{det}}(d_k)\) then renormalize. Peak is
\(\arg\max b_k\). Centroid averages cells with
\(b_k\ge\alpha b_{\max}\) (\(\alpha=0.5\)).

\section{Dipole least-squares fix}
Buffer samples with \(a_i\ge a_{\min}\); fit \((t_x,t_y,A)\) in the same
\(A/(r^{3}+s^{3})\) family. Continuous estimate typically reaches few-metre
error after spiral/verify sampling (\texttt{/swarm/belief/fix}).

\section{Coverage and hunt geometry}
\begin{description}[style=nextline,leftmargin=1.2cm]
  \item[Strip lawnmower (v3 default)] Split \([x_{\min},x_{\max}]\) into \(N\)
        vertical strips with dead gap \(g=20\,\mathrm{m}\); row spacing
        \(15\,\mathrm{m}\). Avoids seam meetings of naive Voronoi partitions.
  \item[Mutual information] Score ring candidates by expected entropy reduction
        under the HIT/MISS model; replan periodically (radii \(10/20/30\,\mathrm{m}\)).
  \item[Expanding spiral (inland-safe)] Concentric rings about the centroid.
        Radius is \emph{shrunk} so the full circle stays inside the coverage box
        with inset \(10\,\mathrm{m}\) --- never clamp points onto the shore wall
        (that piled waypoints on \(x=\pm120\) and threw boats through Gazebo
        colliders).
  \item[Verify orbit] Circle about candidate (radius \(\le\min(20,r_{\mathrm{fit}})\));
        peer entry angle opposite the discoverer; long approaches use \(\sim40\,\mathrm{m}\) transit hops.
\end{description}

\section{LOS guidance (short)}
\[
\psi_{\mathrm{cmd}}=\psi_{\mathrm{path}}+\operatorname{atan2}(-e_{\perp},L),\qquad L=15\,\mathrm{m}.
\]
Empty \texttt{plan} clears the follower (HOLD $\Rightarrow$ zero \texttt{cmd\_vel}).
Thrust mixer: \(T_{L,R}=\mathrm{clip}((v\mp\omega g)S)\).

\chapter{Mission finite-state machine and verify handshake}
\label{ch:fsm}

\section{Modes published on \texttt{mission/state}}
\begin{center}
\begin{tabular}{@{}lll@{}}
\toprule
\texttt{mode} & \texttt{hunt\_phase} & Meaning \\
\midrule
\texttt{GLOBAL\_SEARCH} & (empty) & Regional strip lawnmower \\
\texttt{TARGET\_SEARCH} & \texttt{INFO\_GAIN} & MI surfing \\
\texttt{TARGET\_SEARCH} & \texttt{SPIRAL} & Densify samples (inland-capped) \\
\texttt{HOLD} & (empty) & Discoverer waits / non-verifier clears path \\
\texttt{VERIFY} & (empty) & Orbit + confirmations \\
\texttt{COMPLETE} & (empty) & Done / halted \\
\bottomrule
\end{tabular}
\end{center}

\section{Transition conditions (parameter names)}
\begin{itemize}[leftmargin=*]
  \item \textbf{GLOBAL $\rightarrow$ TARGET:} calibration READY (if required);
        \(t\ge\texttt{min\_seconds\_before\_switch}\) (20\,s); for
        \texttt{consecutive\_high\_p\_required}=4 ticks,
        \(p^\star\ge\texttt{p\_enter\_target\_search}=0.25\) and estimate within
        \texttt{peak\_near\_asv\_radius\_m}=150\,m.
  \item \textbf{INFO\_GAIN $\rightarrow$ SPIRAL:} step budget (45), MI collapse,
        proximity to estimate (\(\le25\,\mathrm{m}\)), or high \(p^\star\) near estimate.
  \item \textbf{SPIRAL $\rightarrow$ HOLD + declare:} after
        \texttt{spiral\_min\_duration\_s}=45\,s with peak/proximity gates (or
        $2\times$ fallback). Publishes empty \texttt{plan} (hard stop).
  \item \textbf{HOLD $\rightarrow$ VERIFY:} matching
        \texttt{/swarm/verify/request} for this \texttt{asv\_id}. Multi-ASV sets
        \texttt{self\_verify: false} so the discoverer does not self-orbit;
        launch forces \texttt{self\_verify:=true} only when \texttt{num\_asvs}=1.
  \item \textbf{Non-verifier $\rightarrow$ HOLD:} any verify request assigned to
        another boat clears our path so we do not pile into the same spiral.
  \item \textbf{VERIFY $\rightarrow$ COMPLETE:} confirmation count reaches
        \texttt{verify\_confirmations\_required}=4
        (at-site \(\le30\,\mathrm{m}\), peak aligned \(\le50\,\mathrm{m}\),
        anomaly \(\ge15\,\mathrm{nT}\), \(p^\star\ge0.30\)).
\end{itemize}

\section{Handshake topics}
\begin{enumerate}[leftmargin=*]
  \item Discoverer $\rightarrow$ \texttt{/swarm/verify/declare}
        (\texttt{VerifyRequest}, verifier\_id empty).
  \item Coordinator $\rightarrow$ \texttt{/swarm/verify/request}
        (fills \texttt{verifier\_id}, prefers other ASV).
  \item Verifier $\rightarrow$ \texttt{/swarm/verify/result}
        (\texttt{VerifyResult}).
  \item Coordinator $\rightarrow$ \texttt{/swarm/mission/complete} and
        \texttt{/swarm/mission/halt} (latched TRANSIENT\_LOCAL).
\end{enumerate}

\section{Plotter status strip}
The trajectory plotter banner mirrors this FSM for operators:
SEARCH / HUNT / VERIFY (with \(n/N\) confirmations) / CONFIRMED.
CONFIRMED is a \emph{software} latch; offline localisation error vs the planted
target is scored separately from \texttt{/swarm/belief/fix}.

\chapter{ROS 2 Node Catalog --- Magnetic ASV Stack (\texttt{version\_3})}
\label{ch:ros-catalog}

This chapter is the authoritative software reference for every message, bridge, and node in the dual-ASV stack. Each node section lists identity, parameters, subscriptions, publications, timers, QoS, internal algorithm steps, and runtime dependencies.


\textbf{Scope:} \texttt{/home/robot/simulation\_ws/version\_3}

\textbf{Packages:} \texttt{boat\_control}, \texttt{boat\_sensing}, \texttt{boat\_mapping}, \texttt{boat\_mission}, \texttt{boat\_navigation}, plus \texttt{boat\_bringup} Gazebo bridges and \texttt{boat\_msgs}.

\textbf{Source of truth:} executable Python nodes, \texttt{setup.py} entry points, \texttt{declare\_parameter} defaults, and \texttt{boat\_bringup/config/*.yaml} overrides as of this tree.


\textbf{Topic naming convention used throughout}


\begin{center}
\small
\begin{longtable}{@{}lp{10cm}@{}}
\toprule
Form & Meaning \\
\midrule
\endfirsthead
\toprule
Form & Meaning \\
\midrule
\endhead
\emph{relative} & No leading \texttt{/} --- resolved under the node's ROS namespace (e.g. node in \texttt{/asv1} → \texttt{cmd\_vel} becomes \texttt{/asv1/cmd\_vel}). \\
\emph{absolute} & Leading \texttt{/} --- global name, not remapped by namespace (e.g. \texttt{/swarm/belief/peak}, \texttt{/model/simple\_boat/...}). \\
\bottomrule
\end{longtable}
\end{center}


\textbf{Typical launch namespaces} (\texttt{sim.launch.py}): ASV1 → \texttt{/asv1} + Gazebo model \texttt{simple\_boat}; ASV2 → \texttt{/asv2} + \texttt{simple\_boat\_2}. Swarm fusion / verify coordinator run at root (no ASV namespace).


\textbf{Data-flow overview (one ASV + swarm)}


\begin{verbatim}
Gazebo ──(gz transport / ros_gz_bridge)──► odom, imu, gps, mag/gazebo, cmd_thrust
gazebo_pose2d ──► pose2d
los_path_follower ◄── odom, plan ──► cmd_vel ──► thrust_mixer ──► cmd_thrust
mag_driver ◄── mag/gazebo?, pose2d, cmd_vel ──► mag/raw
mag_filter ──► mag/filtered
calibration_node ──► mag/anomaly, calibration/status
bayes_fusion ◄── mag/anomaly(+) ──► /swarm/belief/*
mission_manager ◄── pose, belief, anomaly, cal ──► plan, mission/*, verify declare
verify_coordinator ◄── declare, result ──► request, halt, complete
trajectory_plotter (viz)
\end{verbatim}



\chapter{Part I --- Custom Messages (\texttt{boat\_msgs})}
\label{ch:ros-catalog}

This chapter is the authoritative software reference for every message, bridge, and node in the dual-ASV stack. Each node section lists identity, parameters, subscriptions, publications, timers, QoS, internal algorithm steps, and runtime dependencies.


\section{\texttt{MagReading.msg}}


\begin{center}
\small
\begin{longtable}{@{}llp{7.2cm}@{}}
\toprule
Field & Type & Meaning \\
\midrule
\endfirsthead
\toprule
Field & Type & Meaning \\
\midrule
\endhead
\texttt{header} & \texttt{std\_msgs/Header} & Stamp + \texttt{frame\_id} of the magnetometer sample. \\
\texttt{bx}, \texttt{by}, \texttt{bz} & \texttt{float64} & Magnetic field components in nanotesla (nT). \\
\texttt{scalar} & \texttt{float64} & \(\sqrt{b\_x^2+b\_y^2+b\_z^2}\) magnitude (nT). \\
\texttt{x}, \texttt{y} & \texttt{float64} & Planar world pose at sample time (m). \\
\texttt{heading} & \texttt{float64} & Yaw / heading at sample time (rad). \\
\texttt{motor\_on} & \texttt{bool} & True when commanded thruster surge or yaw exceeds the driver's motor threshold (for future motor-state bins). \\
\bottomrule
\end{longtable}
\end{center}


\section{\texttt{MagAnomaly.msg}}


\begin{center}
\small
\begin{longtable}{@{}llp{7.2cm}@{}}
\toprule
Field & Type & Meaning \\
\midrule
\endfirsthead
\toprule
Field & Type & Meaning \\
\midrule
\endhead
\texttt{header} & \texttt{std\_msgs/Header} & Propagated from the filtered \texttt{MagReading}. \\
\texttt{raw\_nt} & \texttt{float64} & Scalar field used as raw input to calibration (nT). \\
\texttt{baseline\_nt} & \texttt{float64} & Estimated ambient baseline (nT); \texttt{nan} if not calibrated. \\
\texttt{cleaned\_anomaly\_nt} & \texttt{float64} & Spatially/temporally cleaned anomaly feature for Bayes / mission (nT). \\
\texttt{is\_calibrated} & \texttt{bool} & True when a usable baseline estimate exists for this sample. \\
\texttt{heading\_bin} & \texttt{int32} & Discrete heading bin used for baseline lookup. \\
\texttt{grid\_cell\_x}, \texttt{grid\_cell\_y} & \texttt{int32} & Baseline-map cell indices \((i,j)\). \\
\texttt{x}, \texttt{y}, \texttt{heading} & \texttt{float64} & Pose/heading of the sample (m, rad). \\
\bottomrule
\end{longtable}
\end{center}


\section{\texttt{CalibrationStatus.msg}}


\begin{center}
\small
\begin{longtable}{@{}p{2.1cm}p{2.1cm}p{2.1cm}p{2.1cm}p{3.2cm}@{}}
\toprule
Field & Type & Meaning &  &  \\
\midrule
\endfirsthead
\toprule
Field & Type & Meaning &  &  \\
\midrule
\endhead
\texttt{header} & \texttt{std\_msgs/Header} & Status stamp; \texttt{frame\_id} set to \texttt{map}. &  &  \\
\texttt{phase} & \texttt{string} & \texttt{CALIBRATING} \ & \texttt{READY} \ & \texttt{DEGRADED}. \\
\texttt{cells\_sampled} & \texttt{int32} & Number of baseline cells that have received samples. &  &  \\
\texttt{coverage\_percent} & \texttt{float32} & Percent of grid cells sampled. &  &  \\
\bottomrule
\end{longtable}
\end{center}


\section{\texttt{BeliefGrid.msg}}


\begin{center}
\small
\begin{longtable}{@{}llp{7.2cm}@{}}
\toprule
Field & Type & Meaning \\
\midrule
\endfirsthead
\toprule
Field & Type & Meaning \\
\midrule
\endhead
\texttt{header} & \texttt{std\_msgs/Header} & Map stamp; \texttt{frame\_id} = \texttt{map}. \\
\texttt{resolution} & \texttt{float32} & Cell size (m). \\
\texttt{origin\_x}, \texttt{origin\_y} & \texttt{float64} & World coordinates of the grid origin (m). \\
\texttt{width}, \texttt{height} & \texttt{uint32} & Grid dimensions in cells. \\
\texttt{data} & \texttt{float64[]} & Row-major posterior probabilities; should sum ≈ 1. \\
\bottomrule
\end{longtable}
\end{center}


\section{\texttt{MissionState.msg}}


\begin{center}
\small
\begin{longtable}{@{}p{2.1cm}p{2.1cm}p{2.1cm}p{2.1cm}p{2.1cm}p{2.1cm}p{3.2cm}@{}}
\toprule
Field & Type & Meaning &  &  &  &  \\
\midrule
\endfirsthead
\toprule
Field & Type & Meaning &  &  &  &  \\
\midrule
\endhead
\texttt{header} & \texttt{std\_msgs/Header} & Snapshot stamp; \texttt{frame\_id} = \texttt{map}. &  &  &  &  \\
\texttt{mode} & \texttt{string} & \texttt{GLOBAL\_SEARCH} \ & \texttt{TARGET\_SEARCH} \ & \texttt{HOLD} \ & \texttt{VERIFY} \ & \texttt{COMPLETE} (and related). \\
\texttt{hunt\_phase} & \texttt{string} & Within target hunt: \texttt{INFO\_GAIN} \ & \texttt{SPIRAL}, else empty. &  &  &  \\
\texttt{peak\_p} & \texttt{float64} & Latest swarm peak belief probability. &  &  &  &  \\
\texttt{peak\_x}, \texttt{peak\_y} & \texttt{float64} & Latest discrete peak cell centre (m). &  &  &  &  \\
\texttt{info\_gain\_steps} & \texttt{int32} & Number of MI replan steps taken in \texttt{INFO\_GAIN}. &  &  &  &  \\
\texttt{confirmations} & \texttt{int32} & Verification confirmation count. &  &  &  &  \\
\bottomrule
\end{longtable}
\end{center}


\section{\texttt{VerifyRequest.msg}}


\begin{center}
\small
\begin{longtable}{@{}llp{7.2cm}@{}}
\toprule
Field & Type & Meaning \\
\midrule
\endfirsthead
\toprule
Field & Type & Meaning \\
\midrule
\endhead
\texttt{header} & \texttt{std\_msgs/Header} & Request stamp; \texttt{frame\_id} = \texttt{map}. \\
\texttt{discoverer\_id} & \texttt{string} & ASV that declared the candidate (e.g. \texttt{asv1}). \\
\texttt{verifier\_id} & \texttt{string} & Assigned verifier; empty on declare, filled by coordinator on request. \\
\texttt{candidate\_x}, \texttt{candidate\_y} & \texttt{float64} & Candidate location (weighted centroid / estimate) (m). \\
\texttt{candidate\_peak\_p} & \texttt{float64} & Peak probability at declaration. \\
\texttt{discoverer\_x}, \texttt{discoverer\_y} & \texttt{float64} & Discoverer pose at declaration (for opposite-side orbit approach). \\
\bottomrule
\end{longtable}
\end{center}


\section{\texttt{VerifyResult.msg}}


\begin{center}
\small
\begin{longtable}{@{}llp{7.2cm}@{}}
\toprule
Field & Type & Meaning \\
\midrule
\endfirsthead
\toprule
Field & Type & Meaning \\
\midrule
\endhead
\texttt{header} & \texttt{std\_msgs/Header} & Result stamp; \texttt{frame\_id} = \texttt{map}. \\
\texttt{success} & \texttt{bool} & True if required confirmations were collected. \\
\texttt{verifier\_id} & \texttt{string} & ASV that performed verification. \\
\texttt{candidate\_x}, \texttt{candidate\_y} & \texttt{float64} & Candidate under verification (m). \\
\texttt{confirmations} & \texttt{int32} & Number of confirming readings. \\
\texttt{final\_peak\_p} & \texttt{float64} & Swarm peak probability at finish. \\
\bottomrule
\end{longtable}
\end{center}




\chapter{Part II --- Gazebo ↔ ROS Bridges (\texttt{boat\_bringup})}
\label{ch:ros-catalog}

This chapter is the authoritative software reference for every message, bridge, and node in the dual-ASV stack. Each node section lists identity, parameters, subscriptions, publications, timers, QoS, internal algorithm steps, and runtime dependencies.


Bridges are \textbf{not} Python package executables; they are \texttt{ros\_gz\_bridge} \texttt{parameter\_bridge} nodes started from \texttt{sim.launch.py}.


\section{Bridge nodes}


\begin{center}
\small
\begin{longtable}{@{}lllp{5.2cm}@{}}
\toprule
Launch name & Executable & Config & When \\
\midrule
\endfirsthead
\toprule
Launch name & Executable & Config & When \\
\midrule
\endhead
\texttt{boat\_bridge} & \texttt{ros\_gz\_bridge/parameter\_bridge} & \texttt{config/bridge.yaml} & Always with \texttt{sim.launch.py} \\
\texttt{clock\_bridge} & \texttt{ros\_gz\_bridge/parameter\_bridge} & \texttt{config/clock\_bridge.yaml} & Only when launch arg \texttt{fast:=true} \\
\bottomrule
\end{longtable}
\end{center}


\textbf{Dependencies:} Gazebo Sim world running (\texttt{gz sim … water\_world.sdf}); thruster topics require \texttt{thrust\_mixer} publishing on the matching \texttt{/model/.../cmd\_thrust} names.


\section{\texttt{bridge.yaml} topic mappings}


\subsection{ASV1 (\texttt{simple\_boat})}


\begin{center}
\small
\begin{longtable}{@{}p{2.1cm}p{2.1cm}p{2.1cm}p{2.1cm}p{3.2cm}@{}}
\toprule
ROS topic & GZ topic & ROS type & GZ type & Direction \\
\midrule
\endfirsthead
\toprule
ROS topic & GZ topic & ROS type & GZ type & Direction \\
\midrule
\endhead
\texttt{/model/simple\_boat/joint/left\_propeller\_joint/cmd\_thrust} & same & \texttt{std\_msgs/msg/Float64} & \texttt{gz.msgs.Double} & ROS\_TO\_GZ \\
\texttt{/model/simple\_boat/joint/right\_propeller\_joint/cmd\_thrust} & same & \texttt{std\_msgs/msg/Float64} & \texttt{gz.msgs.Double} & ROS\_TO\_GZ \\
\texttt{/asv1/imu/data} & \texttt{/model/simple\_boat/imu} & \texttt{sensor\_msgs/msg/Imu} & \texttt{gz.msgs.IMU} & GZ\_TO\_ROS \\
\texttt{/asv1/gps/fix} & \texttt{/model/simple\_boat/gps} & \texttt{sensor\_msgs/msg/NavSatFix} & \texttt{gz.msgs.NavSat} & GZ\_TO\_ROS \\
\texttt{/asv1/mag/gazebo} & \texttt{/model/simple\_boat/magnetometer} & \texttt{sensor\_msgs/msg/MagneticField} & \texttt{gz.msgs.Magnetometer} & GZ\_TO\_ROS \\
\texttt{/asv1/odom} & \texttt{/model/simple\_boat/odometry} & \texttt{nav\_msgs/msg/Odometry} & \texttt{gz.msgs.Odometry} & GZ\_TO\_ROS \\
\bottomrule
\end{longtable}
\end{center}


\subsection{ASV2 (\texttt{simple\_boat\_2})}


\begin{center}
\small
\begin{longtable}{@{}p{2.1cm}p{2.1cm}p{2.1cm}p{2.1cm}p{3.2cm}@{}}
\toprule
ROS topic & GZ topic & ROS type & GZ type & Direction \\
\midrule
\endfirsthead
\toprule
ROS topic & GZ topic & ROS type & GZ type & Direction \\
\midrule
\endhead
\texttt{/model/simple\_boat\_2/joint/left\_propeller\_joint/cmd\_thrust} & same & \texttt{std\_msgs/msg/Float64} & \texttt{gz.msgs.Double} & ROS\_TO\_GZ \\
\texttt{/model/simple\_boat\_2/joint/right\_propeller\_joint/cmd\_thrust} & same & \texttt{std\_msgs/msg/Float64} & \texttt{gz.msgs.Double} & ROS\_TO\_GZ \\
\texttt{/asv2/imu/data} & \texttt{/model/simple\_boat\_2/imu} & \texttt{sensor\_msgs/msg/Imu} & \texttt{gz.msgs.IMU} & GZ\_TO\_ROS \\
\texttt{/asv2/gps/fix} & \texttt{/model/simple\_boat\_2/gps} & \texttt{sensor\_msgs/msg/NavSatFix} & \texttt{gz.msgs.NavSat} & GZ\_TO\_ROS \\
\texttt{/asv2/mag/gazebo} & \texttt{/model/simple\_boat\_2/magnetometer} & \texttt{sensor\_msgs/msg/MagneticField} & \texttt{gz.msgs.Magnetometer} & GZ\_TO\_ROS \\
\texttt{/asv2/odom} & \texttt{/model/simple\_boat\_2/odometry} & \texttt{nav\_msgs/msg/Odometry} & \texttt{gz.msgs.Odometry} & GZ\_TO\_ROS \\
\bottomrule
\end{longtable}
\end{center}


\section{\texttt{clock\_bridge.yaml}}


\begin{center}
\small
\begin{longtable}{@{}p{2.1cm}p{2.1cm}p{2.1cm}p{2.1cm}p{3.2cm}@{}}
\toprule
ROS topic & GZ topic & ROS type & GZ type & Direction \\
\midrule
\endfirsthead
\toprule
ROS topic & GZ topic & ROS type & GZ type & Direction \\
\midrule
\endhead
\texttt{/clock} & \texttt{/clock} & \texttt{rosgraph\_msgs/msg/Clock} & \texttt{gz.msgs.Clock} & GZ\_TO\_ROS \\
\bottomrule
\end{longtable}
\end{center}


\section{Related config (not bridges)}


\texttt{record\_topics.yaml} lists bag topics for offline replay (\texttt{/asv1/pose2d}, mag chain, belief topics, etc.); it does not start a node by itself.




\chapter{Part III --- Nodes by Package}
\label{ch:ros-catalog}

This chapter is the authoritative software reference for every message, bridge, and node in the dual-ASV stack. Each node section lists identity, parameters, subscriptions, publications, timers, QoS, internal algorithm steps, and runtime dependencies.




\section{1. \texttt{boat\_control} --- \texttt{thrust\_mixer}}


\subsection{1.1 Identity}


\begin{center}
\small
\begin{longtable}{@{}lp{10cm}@{}}
\toprule
Item & Value \\
\midrule
\endfirsthead
\toprule
Item & Value \\
\midrule
\endhead
Package & \texttt{boat\_control} \\
Entry point & \texttt{thrust\_mixer = boat\_control.mixer:main} \\
Class & \texttt{ThrustMixer} \\
File & \texttt{src/boat\_control/boat\_control/mixer.py} \\
Core helper & \texttt{boat\_control.core.mix\_thrust} \\
ROS node name & \texttt{thrust\_mixer} \\
\bottomrule
\end{longtable}
\end{center}


\subsection{1.2 Parameters}


\begin{center}
\small
\begin{longtable}{@{}llp{7.2cm}@{}}
\toprule
Name & Code default & YAML / launch override \\
\midrule
\endfirsthead
\toprule
Name & Code default & YAML / launch override \\
\midrule
\endhead
\texttt{model\_name} & \texttt{'simple\_boat'} & Launch: \texttt{simple\_boat} / \texttt{simple\_boat\_2} per ASV. \textbf{No} \texttt{boat\_bringup/config} YAML. \\
\texttt{cmd\_vel\_topic} & \texttt{'cmd\_vel'} & Launch: \texttt{'cmd\_vel'} (relative). \\
\texttt{thrust\_scale} & \texttt{50.0} & Launch: \texttt{50.0} \\
\texttt{turn\_gain} & \texttt{1.0} & Launch: \texttt{1.0} \\
\texttt{max\_thrust} & \texttt{100.0} & Launch: \texttt{100.0} \\
\bottomrule
\end{longtable}
\end{center}


\subsection{1.3 Subscriptions}


\begin{center}
\small
\begin{longtable}{@{}lllp{5.2cm}@{}}
\toprule
Topic & Rel/Abs & Type & Callback purpose \\
\midrule
\endfirsthead
\toprule
Topic & Rel/Abs & Type & Callback purpose \\
\midrule
\endhead
\texttt{cmd\_vel} (param) & relative & \texttt{geometry\_msgs/Twist} & \texttt{on\_cmd\_vel}: convert surge \texttt{linear.x} and yaw rate \texttt{angular.z} into left/right thrust and publish immediately. \\
\bottomrule
\end{longtable}
\end{center}


\subsection{1.4 Publishers}


\begin{center}
\small
\begin{longtable}{@{}lllp{5.2cm}@{}}
\toprule
Topic & Rel/Abs & Type & When \\
\midrule
\endfirsthead
\toprule
Topic & Rel/Abs & Type & When \\
\midrule
\endhead
\texttt{/model/\{model\_name\}/joint/left\_propeller\_joint/cmd\_thrust} & \textbf{absolute} & \texttt{std\_msgs/Float64} & On every \texttt{cmd\_vel} \\
\texttt{/model/\{model\_name\}/joint/right\_propeller\_joint/cmd\_thrust} & \textbf{absolute} & \texttt{std\_msgs/Float64} & On every \texttt{cmd\_vel} \\
\bottomrule
\end{longtable}
\end{center}


\subsection{1.5 Timers / QoS}


\begin{itemize}[leftmargin=*]
\item \textbf{Timers:} none (event-driven).
\item \textbf{QoS:} depth 10, default (reliable, volatile).
\end{itemize}


\subsection{1.6 Algorithm pipeline}


\begin{itemize}[leftmargin=*]
\item Read \texttt{linear.x} (surge) and \texttt{angular.z} (yaw).
\item \(L = (u - r \cdot g) s\), \(R = (u + r \cdot g) s\) with scale \(s\), turn gain \(g\).
\item Clamp each channel to \(\pm\) \texttt{max\_thrust}.
\item Publish left/right \texttt{Float64} thrust commands to Gazebo joints (via bridge).
\end{itemize}


\subsection{1.7 Dependencies}


\begin{itemize}[leftmargin=*]
\item \textbf{Upstream:} \texttt{los\_path\_follower} (or teleop) publishing \texttt{cmd\_vel}.
\item \textbf{Downstream / infra:} \texttt{boat\_bridge} mapping thrust topics into Gazebo; Gazebo model spawned.
\end{itemize}




\section{2. \texttt{boat\_sensing} --- \texttt{mag\_driver}}


\subsection{2.1 Identity}


\begin{center}
\small
\begin{longtable}{@{}lp{10cm}@{}}
\toprule
Item & Value \\
\midrule
\endfirsthead
\toprule
Item & Value \\
\midrule
\endhead
Entry point & \texttt{mag\_driver = boat\_sensing.mag\_driver:main} \\
Class & \texttt{MagDriver} \\
File & \texttt{src/boat\_sensing/boat\_sensing/mag\_driver.py} \\
\bottomrule
\end{longtable}
\end{center}


\subsection{2.2 Parameters}


\begin{center}
\small
\begin{longtable}{@{}llp{7.2cm}@{}}
\toprule
Name & Code default & \texttt{sensing.yaml} (\texttt{/**/mag\_driver}) \\
\midrule
\endfirsthead
\toprule
Name & Code default & \texttt{sensing.yaml} (\texttt{/**/mag\_driver}) \\
\midrule
\endhead
\texttt{publish\_rate\_hz} & \texttt{20.0} & \texttt{20.0} \\
\texttt{gazebo\_mag\_topic} & \texttt{'mag/gazebo'} & \texttt{mag/gazebo} \\
\texttt{pose\_topic} & \texttt{'pose2d'} & \texttt{pose2d} \\
\texttt{cmd\_vel\_topic} & \texttt{'cmd\_vel'} & \texttt{cmd\_vel} \\
\texttt{raw\_topic} & \texttt{'mag/raw'} & \texttt{mag/raw} \\
\texttt{frame\_id} & \texttt{'asv1/mag\_link'} & Launch overrides to \texttt{\{ns\}/mag\_link} \\
\texttt{motor\_on\_threshold} & \texttt{1.0e-3} & \texttt{0.001} \\
\texttt{plant\_magnetic\_target} & \texttt{True} & \texttt{true} \\
\texttt{target\_x} & \texttt{80.0} & \textbf{\texttt{55.0}} \\
\texttt{target\_y} & \texttt{-40.0} & \textbf{\texttt{-35.0}} \\
\texttt{dipole\_strength\_nt} & \texttt{1.5e12} & \textbf{\texttt{4.0e5}} \\
\texttt{dipole\_soft\_m} & \texttt{1.0} & \textbf{\texttt{20.0}} \\
\texttt{synthetic\_background\_nt} & \texttt{4.5e8} & \textbf{\texttt{45000.0}} \\
\texttt{synthetic\_noise\_nt} & \texttt{5.0e4} & \textbf{\texttt{3.0}} \\
\bottomrule
\end{longtable}
\end{center}


\subsection{2.3 Subscriptions}


\begin{center}
\small
\begin{longtable}{@{}p{2.1cm}p{2.1cm}p{2.1cm}p{2.1cm}p{2.1cm}p{2.1cm}p{2.1cm}p{3.2cm}@{}}
\toprule
Topic & Rel/Abs & Type & Callback purpose &  &  &  &  \\
\midrule
\endfirsthead
\toprule
Topic & Rel/Abs & Type & Callback purpose &  &  &  &  \\
\midrule
\endhead
\texttt{mag/gazebo} & relative & \texttt{sensor\_msgs/MagneticField} & Cache latest Gazebo magnetometer (used when planting disabled). &  &  &  &  \\
\texttt{pose2d} & relative & \texttt{geometry\_msgs/Pose2D} & Cache pose for stamping / synthetic dipole. &  &  &  &  \\
\texttt{cmd\_vel} & relative & \texttt{geometry\_msgs/Twist} & Set \texttt{motor\_on} if \ & surge\ & or \ & yaw\ & exceeds threshold. \\
\bottomrule
\end{longtable}
\end{center}


\subsection{2.4 Publishers}


\begin{center}
\small
\begin{longtable}{@{}lllp{5.2cm}@{}}
\toprule
Topic & Rel/Abs & Type & When \\
\midrule
\endfirsthead
\toprule
Topic & Rel/Abs & Type & When \\
\midrule
\endhead
\texttt{mag/raw} & relative & \texttt{boat\_msgs/MagReading} & Timer at \texttt{publish\_rate\_hz} (skip if planted mode and no pose, or non-planted and no mag). \\
\bottomrule
\end{longtable}
\end{center}


\subsection{2.5 Timers / QoS}


\begin{itemize}[leftmargin=*]
\item \textbf{Timer:} \texttt{1/publish\_rate\_hz} (default \textbf{20 Hz}) → \texttt{on\_timer}.
\item \textbf{QoS:} mag sub depth 50; others 10; default QoS.
\end{itemize}


\subsection{2.6 Algorithm pipeline}


\begin{itemize}[leftmargin=*]
\item If \texttt{plant\_magnetic\_target}: require pose; compute \(A/(r^3+s^3)\) dipole anomaly; set \(b\_z =\) background + anomaly + Gaussian noise; \(b\_x=b\_y=0\).
\item Else: convert Gazebo Tesla → nT; attach pose if available (else NaN).
\item Fill \texttt{MagReading} (axes, scalar, motor\_on, pose, heading); publish.
\end{itemize}


\subsection{2.7 Dependencies}


\begin{itemize}[leftmargin=*]
\item \textbf{Required for planted mode:} \texttt{gazebo\_pose2d} (\texttt{pose2d}).
\item \textbf{Optional:} Gazebo mag via bridge (\texttt{mag/gazebo}) when planting off; \texttt{cmd\_vel} for \texttt{motor\_on}.
\item \textbf{Consumers:} \texttt{mag\_filter}.
\end{itemize}




\section{3. \texttt{boat\_sensing} --- \texttt{mag\_filter}}


\subsection{3.1 Identity}


\begin{center}
\small
\begin{longtable}{@{}lp{10cm}@{}}
\toprule
Item & Value \\
\midrule
\endfirsthead
\toprule
Item & Value \\
\midrule
\endhead
Entry point & \texttt{mag\_filter = boat\_sensing.mag\_filter:main} \\
Class & \texttt{MagFilter} \\
File & \texttt{src/boat\_sensing/boat\_sensing/mag\_filter.py} \\
Core & \texttt{MagnetometerFilterChain} in \texttt{filter\_core.py} \\
\bottomrule
\end{longtable}
\end{center}


\subsection{3.2 Parameters}


\begin{center}
\small
\begin{longtable}{@{}llp{7.2cm}@{}}
\toprule
Name & Code default & \texttt{sensing.yaml} (\texttt{/**/mag\_filter}) \\
\midrule
\endfirsthead
\toprule
Name & Code default & \texttt{sensing.yaml} (\texttt{/**/mag\_filter}) \\
\midrule
\endhead
\texttt{raw\_topic} & \texttt{'mag/raw'} & same \\
\texttt{filtered\_topic} & \texttt{'mag/filtered'} & same \\
\texttt{status\_topic} & \texttt{'mag/filter\_status'} & same \\
\texttt{lowpass\_window} & \texttt{5} & \texttt{5} \\
\texttt{spike\_history} & \texttt{20} & \texttt{20} \\
\texttt{spike\_n\_sigma} & \texttt{3.0} & \textbf{\texttt{10.0}} \\
\texttt{min\_std\_nt} & \texttt{1.0} & \texttt{1.0} \\
\texttt{status\_rate\_hz} & \texttt{1.0} & \texttt{1.0} \\
\bottomrule
\end{longtable}
\end{center}


\subsection{3.3 Subscriptions}


\begin{center}
\small
\begin{longtable}{@{}lllp{5.2cm}@{}}
\toprule
Topic & Rel/Abs & Type & Callback purpose \\
\midrule
\endfirsthead
\toprule
Topic & Rel/Abs & Type & Callback purpose \\
\midrule
\endhead
\texttt{mag/raw} & relative & \texttt{boat\_msgs/MagReading} & Spike-reject then low-pass each axis/scalar; republish cleaned reading with pose fields copied. \\
\bottomrule
\end{longtable}
\end{center}


\subsection{3.4 Publishers}


\begin{center}
\small
\begin{longtable}{@{}lllp{5.2cm}@{}}
\toprule
Topic & Rel/Abs & Type & When \\
\midrule
\endfirsthead
\toprule
Topic & Rel/Abs & Type & When \\
\midrule
\endhead
\texttt{mag/filtered} & relative & \texttt{boat\_msgs/MagReading} & On every accepted/filtered raw sample \\
\texttt{mag/filter\_status} & relative & \texttt{diagnostic\_msgs/DiagnosticStatus} & Timer at \texttt{status\_rate\_hz} \\
\bottomrule
\end{longtable}
\end{center}


\subsection{3.5 Timers / QoS}


\begin{itemize}[leftmargin=*]
\item \textbf{Timer:} status at \textbf{1 Hz} (default).
\item \textbf{QoS:} raw sub depth 50; pubs depth 10; default.
\end{itemize}


\subsection{3.6 Algorithm pipeline}


\begin{itemize}[leftmargin=*]
\item Per axis + scalar: spike gate vs rolling mean ± \texttt{n\_sigma}·std (min std floor); after 5 consecutive rejects, accept step change.
\item Moving-average low-pass of accepted (or held) value.
\item Copy pose / \texttt{motor\_on} / header; publish; status reports counts and rejects.
\end{itemize}


\subsection{3.7 Dependencies}


\begin{itemize}[leftmargin=*]
\item \textbf{Upstream:} \texttt{mag\_driver}.
\item \textbf{Downstream:} \texttt{calibration\_node}.
\end{itemize}




\section{4. \texttt{boat\_sensing} --- \texttt{calibration\_node}}


\subsection{4.1 Identity}


\begin{center}
\small
\begin{longtable}{@{}lp{10cm}@{}}
\toprule
Item & Value \\
\midrule
\endfirsthead
\toprule
Item & Value \\
\midrule
\endhead
Entry point & \texttt{calibration\_node = boat\_sensing.calibration\_node:main} \\
Class & \texttt{CalibrationNode} \\
File & \texttt{src/boat\_sensing/boat\_sensing/calibration\_node.py} \\
Core & \texttt{MagneticCalibrator}, \texttt{BaselineMap}, \texttt{TemporalHighPass} \\
\bottomrule
\end{longtable}
\end{center}


\subsection{4.2 Parameters}


\begin{center}
\small
\begin{longtable}{@{}llp{7.2cm}@{}}
\toprule
Name & Code default & \texttt{sensing.yaml} (\texttt{/**/calibration\_node}) \\
\midrule
\endfirsthead
\toprule
Name & Code default & \texttt{sensing.yaml} (\texttt{/**/calibration\_node}) \\
\midrule
\endhead
\texttt{filtered\_topic} & \texttt{'mag/filtered'} & same \\
\texttt{anomaly\_topic} & \texttt{'mag/anomaly'} & same \\
\texttt{status\_topic} & \texttt{'calibration/status'} & same \\
\texttt{area\_size\_m} & \texttt{300.0} & \texttt{300.0} \\
\texttt{origin\_x} & \texttt{-150.0} & \texttt{-150.0} \\
\texttt{origin\_y} & \texttt{-150.0} & \texttt{-150.0} \\
\texttt{cell\_size\_m} & \texttt{20.0} & \textbf{\texttt{10.0}} \\
\texttt{num\_heading\_bins} & \texttt{8} & \texttt{8} \\
\texttt{min\_cell\_samples} & \texttt{1} & \texttt{1} \\
\texttt{reject\_residual\_nt} & \texttt{5.0e7} & \textbf{\texttt{35.0}} \\
\texttt{temporal\_window} & \texttt{12} & \texttt{12} \\
\texttt{noise\_floor\_nt} & \texttt{0.0} & \texttt{0.0} \\
\texttt{ready\_coverage\_percent} & \texttt{5.0} & \texttt{5.0} \\
\texttt{ready\_min\_cells} & \texttt{8} & \texttt{8} \\
\texttt{freeze\_baseline\_when\_ready} & \texttt{True} & \texttt{true} \\
\texttt{status\_rate\_hz} & \texttt{1.0} & \texttt{1.0} \\
\bottomrule
\end{longtable}
\end{center}


\subsection{4.3 Subscriptions}


\begin{center}
\small
\begin{longtable}{@{}lllp{5.2cm}@{}}
\toprule
Topic & Rel/Abs & Type & Callback purpose \\
\midrule
\endfirsthead
\toprule
Topic & Rel/Abs & Type & Callback purpose \\
\midrule
\endhead
\texttt{mag/filtered} & relative & \texttt{boat\_msgs/MagReading} & Drop NaN pose; run calibrator; publish \texttt{MagAnomaly}. \\
\bottomrule
\end{longtable}
\end{center}


\subsection{4.4 Publishers}


\begin{center}
\small
\begin{longtable}{@{}lllp{5.2cm}@{}}
\toprule
Topic & Rel/Abs & Type & When \\
\midrule
\endfirsthead
\toprule
Topic & Rel/Abs & Type & When \\
\midrule
\endhead
\texttt{mag/anomaly} & relative & \texttt{boat\_msgs/MagAnomaly} & Per valid filtered sample \\
\texttt{calibration/status} & relative & \texttt{boat\_msgs/CalibrationStatus} & Timer at \texttt{status\_rate\_hz} \\
\bottomrule
\end{longtable}
\end{center}


\subsection{4.5 Timers / QoS}


\begin{itemize}[leftmargin=*]
\item \textbf{Timer:} status \textbf{1 Hz}.
\item \textbf{QoS:} filtered sub depth 50; default elsewhere.
\end{itemize}


\subsection{4.6 Algorithm pipeline}


\begin{itemize}[leftmargin=*]
\item Update (or freeze) spatial baseline map keyed by cell × heading bin; reject learning when residual too large (near dipole).
\item Estimate baseline (own bin → cell mean → neighbour median → global mean).
\item Spatial anomaly = \|raw − baseline\|; temporal high-pass / track boost → \texttt{cleaned\_anomaly\_nt}.
\item Phase → \texttt{READY} when coverage or cell count thresholds met; optionally freeze further baseline learning.
\item Publish anomaly + periodic status.
\end{itemize}


\subsection{4.7 Dependencies}


\begin{itemize}[leftmargin=*]
\item \textbf{Upstream:} \texttt{mag\_filter} (and thus \texttt{mag\_driver} + pose).
\item \textbf{Downstream:} \texttt{bayes\_fusion}, \texttt{mission\_manager}, \texttt{trajectory\_plotter}.
\end{itemize}




\section{5. \texttt{boat\_mapping} --- \texttt{bayes\_fusion}}


\subsection{5.1 Identity}


\begin{center}
\small
\begin{longtable}{@{}lp{10cm}@{}}
\toprule
Item & Value \\
\midrule
\endfirsthead
\toprule
Item & Value \\
\midrule
\endhead
Entry point & \texttt{bayes\_fusion = boat\_mapping.bayes\_fusion:main} \\
Class & \texttt{BayesFusionNode} \\
File & \texttt{src/boat\_mapping/boat\_mapping/bayes\_fusion.py} \\
Core & \texttt{BeliefMap} (\texttt{bayes\_core.py}), \texttt{DipoleFitter} (\texttt{dipole\_fit.py}) \\
Namespace & typically \textbf{root} (not under \texttt{/asv*}) \\
\bottomrule
\end{longtable}
\end{center}


\subsection{5.2 Parameters}


\begin{center}
\small
\begin{longtable}{@{}llp{7.2cm}@{}}
\toprule
Name & Code default & \texttt{mapping.yaml} (\texttt{bayes\_fusion}) \\
\midrule
\endfirsthead
\toprule
Name & Code default & \texttt{mapping.yaml} (\texttt{bayes\_fusion}) \\
\midrule
\endhead
\texttt{anomaly\_topics} & \texttt{['/asv1/mag/anomaly']} & \textbf{\texttt{['/asv1/mag/anomaly','/asv2/mag/anomaly']}} \\
\texttt{map\_topic} & \texttt{/swarm/belief/map} & same \\
\texttt{peak\_topic} & \texttt{/swarm/belief/peak} & same \\
\texttt{peak\_probability\_topic} & \texttt{/swarm/belief/peak\_probability} & same \\
\texttt{centroid\_topic} & \texttt{/swarm/belief/centroid} & same \\
\texttt{centroid\_mass\_topic} & \texttt{/swarm/belief/centroid\_mass} & same \\
\texttt{centroid\_spread\_topic} & \texttt{/swarm/belief/centroid\_spread} & same \\
\texttt{fix\_topic} & \texttt{/swarm/belief/fix} & same \\
\texttt{fix\_rms\_topic} & \texttt{/swarm/belief/fix\_rms} & same \\
\texttt{fix\_samples\_topic} & \texttt{/swarm/belief/fix\_samples} & same \\
\texttt{publish\_rate\_hz} & \texttt{2.0} & \texttt{2.0} \\
\texttt{area\_size\_m} & \texttt{300.0} & \texttt{300.0} \\
\texttt{origin\_x} / \texttt{origin\_y} & \texttt{-150.0} & same \\
\texttt{cell\_size\_m} & \texttt{20.0} & \textbf{\texttt{10.0}} \\
\texttt{p\_bg} & \texttt{0.05} & \texttt{0.05} \\
\texttt{p\_max} & \texttt{0.95} & \texttt{0.95} \\
\texttt{d\_half} & \texttt{30.0} & \texttt{30.0} \\
\texttt{hit\_threshold\_nt} & \texttt{5.0e7} & \textbf{\texttt{15.0}} \\
\texttt{miss\_threshold\_nt} & \texttt{1.0e6} & \textbf{\texttt{5.0}} \\
\texttt{hit\_only} & \texttt{True} & \texttt{true} \\
\texttt{centroid\_threshold\_frac} & \texttt{0.5} & \texttt{0.5} \\
\texttt{dipole\_fit\_enable} & \texttt{True} & \texttt{true} \\
\texttt{dipole\_soft\_m} & \texttt{20.0} & \texttt{20.0} \\
\texttt{dipole\_fit\_min\_anomaly\_nt} & \texttt{10.0} & \texttt{10.0} \\
\texttt{dipole\_fit\_min\_samples} & \texttt{12} & \texttt{12} \\
\texttt{dipole\_fit\_max\_samples} & \texttt{400} & \texttt{400} \\
\texttt{dipole\_fit\_guess\_strength\_nt} & \texttt{4.0e5} & \texttt{4.0e5} \\
\bottomrule
\end{longtable}
\end{center}


\subsection{5.3 Subscriptions}


\begin{center}
\small
\begin{longtable}{@{}lllp{5.2cm}@{}}
\toprule
Topic & Rel/Abs & Type & Callback purpose \\
\midrule
\endfirsthead
\toprule
Topic & Rel/Abs & Type & Callback purpose \\
\midrule
\endhead
each of \texttt{anomaly\_topics} & absolute list & \texttt{boat\_msgs/MagAnomaly} & Bayesian HIT/MISS update; buffer strong samples for dipole LS fit; log HITs. \\
\bottomrule
\end{longtable}
\end{center}


\subsection{5.4 Publishers}


\begin{center}
\small
\begin{longtable}{@{}lllp{5.2cm}@{}}
\toprule
Topic & Rel/Abs & Type & When \\
\midrule
\endfirsthead
\toprule
Topic & Rel/Abs & Type & When \\
\midrule
\endhead
\texttt{/swarm/belief/map} & absolute & \texttt{boat\_msgs/BeliefGrid} & Timer \texttt{publish\_rate\_hz} \\
\texttt{/swarm/belief/peak} & absolute & \texttt{geometry\_msgs/PoseStamped} & same \\
\texttt{/swarm/belief/peak\_probability} & absolute & \texttt{std\_msgs/Float64} & same \\
\texttt{/swarm/belief/centroid} & absolute & \texttt{geometry\_msgs/PoseStamped} & same \\
\texttt{/swarm/belief/centroid\_mass} & absolute & \texttt{std\_msgs/Float64} & same \\
\texttt{/swarm/belief/centroid\_spread} & absolute & \texttt{std\_msgs/Float64} & same \\
\texttt{/swarm/belief/fix} & absolute & \texttt{geometry\_msgs/PoseStamped} & Timer, only if last fit succeeded \\
\texttt{/swarm/belief/fix\_rms} & absolute & \texttt{std\_msgs/Float64} & with fix \\
\texttt{/swarm/belief/fix\_samples} & absolute & \texttt{std\_msgs/Float64} & with fix (sample count as float) \\
\bottomrule
\end{longtable}
\end{center}


\subsection{5.5 Timers / QoS}


\begin{itemize}[leftmargin=*]
\item \textbf{Timer:} belief publish at \textbf{2 Hz} (default/yaml).
\item \textbf{QoS:} anomaly sub depth 50; pubs depth 10; \textbf{no} transient\_local.
\end{itemize}


\subsection{5.6 Algorithm pipeline}


\begin{itemize}[leftmargin=*]
\item Skip uncalibrated samples (ABSTAIN).
\item Classify anomaly: HIT if ≥ hit threshold; MISS if ≤ miss (ignored when \texttt{hit\_only}); else ABSTAIN.
\item For HIT: multiply each cell belief by \(P\_{\mathrm{det}}(d)\) with soft \(1/d^3\) detection model; renormalize.
\item Maintain weighted centroid of cells with belief ≥ \texttt{centroid\_threshold\_frac} × peak.
\item Dipole fitter: buffer strong anomalies; LM fit \(A/(r^3+s^3)\) seeded from centroid → continuous fix.
\item Periodically publish grid, peak, centroid, and optional fix.
\end{itemize}


\subsection{5.7 Dependencies}


\begin{itemize}[leftmargin=*]
\item \textbf{Upstream:} one or more \texttt{calibration\_node} anomaly streams.
\item \textbf{Downstream:} \texttt{mission\_manager}, \texttt{trajectory\_plotter}.
\item Launch condition: \texttt{mapping:=true} in \texttt{sim.launch.py}.
\end{itemize}




\section{6. \texttt{boat\_mission} --- \texttt{mission\_manager}}


\subsection{6.1 Identity}


\begin{center}
\small
\begin{longtable}{@{}lp{10cm}@{}}
\toprule
Item & Value \\
\midrule
\endfirsthead
\toprule
Item & Value \\
\midrule
\endhead
Entry point & \texttt{mission\_manager = boat\_mission.mission\_manager:main} \\
Class & \texttt{MissionManager} \\
File & \texttt{src/boat\_mission/boat\_mission/mission\_manager.py} \\
Helpers & \texttt{InfoGainPlanner}, path generators, \texttt{VerificationTracker} \\
\bottomrule
\end{longtable}
\end{center}


\subsection{6.2 Parameters}


Topic / ID parameters (code defaults; yaml under \texttt{/**/mission\_manager} matches unless noted):


\begin{center}
\small
\begin{longtable}{@{}llp{7.2cm}@{}}
\toprule
Name & Code default & YAML notes \\
\midrule
\endfirsthead
\toprule
Name & Code default & YAML notes \\
\midrule
\endhead
\texttt{asv\_id} & \texttt{''} (→ namespace or \texttt{asv1}) & \texttt{''} \\
\texttt{pose\_topic} & \texttt{pose2d} & same \\
\texttt{plan\_topic} & \texttt{plan} & same \\
\texttt{state\_topic} & \texttt{mission/state} & same \\
\texttt{mode\_topic} & \texttt{mission/mode} & same \\
\texttt{info\_gain\_topic} & \texttt{mission/info\_gain} & same \\
\texttt{anomaly\_topic} & \texttt{mag/anomaly} & same \\
\texttt{calibration\_status\_topic} & \texttt{calibration/status} & same \\
\texttt{peak\_topic} & \texttt{/swarm/belief/peak} & same \\
\texttt{peak\_probability\_topic} & \texttt{/swarm/belief/peak\_probability} & same \\
\texttt{centroid\_topic} & \texttt{/swarm/belief/centroid} & same \\
\texttt{belief\_map\_topic} & \texttt{/swarm/belief/map} & same \\
\texttt{declare\_topic} & \texttt{/swarm/verify/declare} & same \\
\texttt{verify\_request\_topic} & \texttt{/swarm/verify/request} & same \\
\texttt{verify\_result\_topic} & \texttt{/swarm/verify/result} & same \\
\texttt{halt\_topic} & \texttt{/swarm/mission/halt} & same \\
\bottomrule
\end{longtable}
\end{center}


Behaviour / geometry (highlighting yaml overrides):


\begin{center}
\small
\begin{longtable}{@{}llp{7.2cm}@{}}
\toprule
Name & Code default & YAML \\
\midrule
\endfirsthead
\toprule
Name & Code default & YAML \\
\midrule
\endhead
\texttt{min\_x}…\texttt{max\_y} & ±120 & same \\
\texttt{lawnmower\_spacing} & \texttt{40.0} & \textbf{\texttt{20.0}} \\
\texttt{lawnmower\_asv\_index} & \texttt{0} & launch sets per ASV \\
\texttt{lawnmower\_num\_asvs} & \texttt{1} & launch sets to \texttt{num\_asvs} \\
\texttt{lawnmower\_partition} & \texttt{'voronoi'} & \texttt{voronoi} \\
\texttt{region\_seeds} & \texttt{[-110,-110,110,-110]} & same (cleared if \texttt{num\_asvs≤1}) \\
\texttt{control\_rate\_hz} & \texttt{2.0} & \texttt{2.0} \\
\texttt{p\_enter\_target\_search} & \texttt{0.25} & same \\
\texttt{consecutive\_high\_p\_required} & \texttt{4} & same \\
\texttt{peak\_near\_asv\_radius\_m} & \texttt{150.0} & same \\
\texttt{min\_seconds\_before\_switch} & \texttt{20.0} & same \\
\texttt{require\_calibration\_ready} & \texttt{True} & same \\
\texttt{force\_target\_search} / \texttt{force\_spiral\_complete} / \texttt{force\_verify} & \texttt{False} & same \\
\texttt{info\_gain\_*} & see code & yaml mirrors \\
\texttt{spiral\_*} & see code & yaml mirrors \\
\texttt{self\_verify} & \texttt{True} & launch: \texttt{True} if 1 ASV, \texttt{False} if 2 \\
\texttt{verify\_orbit\_radius\_m} & \texttt{20.0} & same \\
\texttt{verify\_orbit\_points} & \texttt{12} & same \\
\texttt{verify\_confirmations\_required} & \texttt{4} & same \\
\texttt{verify\_arrival\_radius\_m} & \texttt{30.0} & same \\
\texttt{verify\_peak\_tolerance\_m} & \texttt{50.0} & same \\
\texttt{verify\_confirmation\_threshold\_nt} & \texttt{5.0e7} & \textbf{\texttt{15.0}} \\
\texttt{verify\_min\_peak\_probability} & \texttt{0.30} & same \\
\bottomrule
\end{longtable}
\end{center}


\subsection{6.3 Subscriptions}


\begin{center}
\small
\begin{longtable}{@{}lllp{5.2cm}@{}}
\toprule
Topic & Rel/Abs & Type & Callback purpose \\
\midrule
\endfirsthead
\toprule
Topic & Rel/Abs & Type & Callback purpose \\
\midrule
\endhead
\texttt{pose2d} & relative & \texttt{Pose2D} & Own pose for planning / distance checks. \\
\texttt{/swarm/belief/peak} & absolute & \texttt{PoseStamped} & Discrete peak cell location. \\
\texttt{/swarm/belief/peak\_probability} & absolute & \texttt{Float64} & Peak probability for mode switch / verify. \\
\texttt{/swarm/belief/centroid} & absolute & \texttt{PoseStamped} & Preferred target estimate (else peak). \\
\texttt{/swarm/belief/map} & absolute & \texttt{BeliefGrid} & Feed MI planner belief. \\
\texttt{mag/anomaly} & relative & \texttt{MagAnomaly} & Latest anomaly; during VERIFY, register confirmations. \\
\texttt{calibration/status} & relative & \texttt{CalibrationStatus} & Gate GLOBAL→TARGET on \texttt{READY}. \\
\texttt{/swarm/verify/request} & absolute & \texttt{VerifyRequest} & If addressed to this ASV, start VERIFY orbit. \\
\texttt{/swarm/mission/halt} & absolute & \texttt{Bool} & On true, enter COMPLETE. \\
\bottomrule
\end{longtable}
\end{center}


\subsection{6.4 Publishers}


\begin{center}
\small
\begin{longtable}{@{}lllp{5.2cm}@{}}
\toprule
Topic & Rel/Abs & Type & When \\
\midrule
\endfirsthead
\toprule
Topic & Rel/Abs & Type & When \\
\midrule
\endhead
\texttt{plan} & relative & \texttt{nav\_msgs/Path} & Lawnmower / MI waypoint / spiral / orbit / hold; \textbf{TRANSIENT\_LOCAL} \\
\texttt{mission/state} & relative & \texttt{MissionState} & Every control tick \\
\texttt{mission/mode} & relative & \texttt{std\_msgs/String} & Every control tick \\
\texttt{mission/info\_gain} & relative & \texttt{Float64} & Every control tick (last MI gain) \\
\texttt{/swarm/verify/declare} & absolute & \texttt{VerifyRequest} & Once after spiral → HOLD \\
\texttt{/swarm/verify/result} & absolute & \texttt{VerifyResult} & End of VERIFY; \textbf{TRANSIENT\_LOCAL} \\
\bottomrule
\end{longtable}
\end{center}


\subsection{6.5 Timers / QoS}


\begin{itemize}[leftmargin=*]
\item \textbf{Timer:} \texttt{control\_rate\_hz} → \textbf{2 Hz} state machine.
\item \textbf{QoS:} \texttt{plan} and \texttt{verify\_result}: Reliable + \textbf{TRANSIENT\_LOCAL}, depth 1. Others depth 10/50 default.
\end{itemize}


\subsection{6.6 Algorithm pipeline}


\begin{itemize}[leftmargin=*]
\item \textbf{GLOBAL\_SEARCH:} publish partitioned lawnmower (voronoi regions or interleaved lanes); wait for cal ready, time, and consecutive high peak near ASV → TARGET\_SEARCH.
\item \textbf{TARGET\_SEARCH / INFO\_GAIN:} mutual-information waypoint selection on belief grid; replan every \texttt{info\_gain\_replan\_period\_s}; exit on max steps / low MI / near peak / confidence.
\item \textbf{SPIRAL:} expanding spiral about centroid/peak; complete on duration + peak\_p + proximity (or force); → HOLD + declare candidate.
\item \textbf{HOLD:} wait for coordinator assignment unless \texttt{self\_verify} starts VERIFY immediately.
\item \textbf{VERIFY:} opposite-side orbit; count confirming anomalies (at site, peak aligned, strong, confident); publish result → COMPLETE.
\item \textbf{COMPLETE / halt:} hold-in-place plan.
\end{itemize}


\subsection{6.7 Dependencies}


\begin{itemize}[leftmargin=*]
\item \texttt{gazebo\_pose2d}, \texttt{los\_path\_follower} (consumes \texttt{plan}), sensing chain + \texttt{bayes\_fusion}, optionally \texttt{verify\_coordinator}.
\item Launch: \texttt{mission:=true} and \texttt{autonomy:=true} for LOS with \texttt{use\_builtin\_path:=false}.
\end{itemize}




\section{7. \texttt{boat\_mission} --- \texttt{verify\_coordinator}}


\subsection{7.1 Identity}


\begin{center}
\small
\begin{longtable}{@{}lp{10cm}@{}}
\toprule
Item & Value \\
\midrule
\endfirsthead
\toprule
Item & Value \\
\midrule
\endhead
Entry point & \texttt{verify\_coordinator = boat\_mission.verify\_coordinator:main} \\
Class & \texttt{VerifyCoordinator} \\
File & \texttt{src/boat\_mission/boat\_mission/verify\_coordinator.py} \\
Namespace & root (launched once, typically with ASV1 stack) \\
\bottomrule
\end{longtable}
\end{center}


\subsection{7.2 Parameters}


\begin{center}
\small
\begin{longtable}{@{}llp{7.2cm}@{}}
\toprule
Name & Code default & \texttt{mission.yaml} (\texttt{verify\_coordinator}) \\
\midrule
\endfirsthead
\toprule
Name & Code default & \texttt{mission.yaml} (\texttt{verify\_coordinator}) \\
\midrule
\endhead
\texttt{declare\_topic} & \texttt{/swarm/verify/declare} & same \\
\texttt{request\_topic} & \texttt{/swarm/verify/request} & same \\
\texttt{result\_topic} & \texttt{/swarm/verify/result} & same \\
\texttt{complete\_topic} & \texttt{/swarm/mission/complete} & same \\
\texttt{halt\_topic} & \texttt{/swarm/mission/halt} & same \\
\texttt{known\_asvs} & \texttt{['asv1']} & \textbf{\texttt{['asv1','asv2']}} \\
\texttt{prefer\_other\_verifier} & \texttt{True} & \texttt{true} \\
\bottomrule
\end{longtable}
\end{center}


\subsection{7.3 Subscriptions}


\begin{center}
\small
\begin{longtable}{@{}lllp{5.2cm}@{}}
\toprule
Topic & Rel/Abs & Type & Callback purpose \\
\midrule
\endfirsthead
\toprule
Topic & Rel/Abs & Type & Callback purpose \\
\midrule
\endhead
\texttt{/swarm/verify/declare} & absolute & \texttt{VerifyRequest} & Pick verifier (prefer other ASV); publish assigned request. \\
\texttt{/swarm/verify/result} & absolute & \texttt{VerifyResult} & On success: latch complete+halt; on failure: clear active request. \\
\bottomrule
\end{longtable}
\end{center}


\subsection{7.4 Publishers}


\begin{center}
\small
\begin{longtable}{@{}lllp{5.2cm}@{}}
\toprule
Topic & Rel/Abs & Type & When \\
\midrule
\endfirsthead
\toprule
Topic & Rel/Abs & Type & When \\
\midrule
\endhead
\texttt{/swarm/verify/request} & absolute & \texttt{VerifyRequest} & On accept declare \\
\texttt{/swarm/mission/complete} & absolute & \texttt{Bool} & Success; \textbf{TRANSIENT\_LOCAL} \\
\texttt{/swarm/mission/halt} & absolute & \texttt{Bool} & Success; \textbf{TRANSIENT\_LOCAL} \\
\texttt{/swarm/mission/status} & absolute (\textbf{hardcoded}) & \texttt{String} & Assignment / complete strings; \textbf{TRANSIENT\_LOCAL} \\
\bottomrule
\end{longtable}
\end{center}


\subsection{7.5 Timers / QoS}


\begin{itemize}[leftmargin=*]
\item \textbf{Timers:} none.
\item \textbf{QoS:} complete / halt / status: Reliable + \textbf{TRANSIENT\_LOCAL}, depth 1. Request/declare: depth 10 default.
\end{itemize}


\subsection{7.6 Algorithm pipeline}


\begin{itemize}[leftmargin=*]
\item Ignore declares while busy or mission already complete.
\item Choose verifier ≠ discoverer if possible from \texttt{known\_asvs}; else self.
\item Forward candidate + discoverer pose as \texttt{VerifyRequest}.
\item On successful \texttt{VerifyResult}: publish \texttt{complete=true}, \texttt{halt=true}, status string.
\end{itemize}


\subsection{7.7 Dependencies}


\begin{itemize}[leftmargin=*]
\item At least one \texttt{mission\_manager} publishing declares / results.
\item Peer \texttt{mission\_manager} instances listening for requests / halt.
\end{itemize}




\section{8. \texttt{boat\_mission} --- \texttt{spiral\_demo}}


\subsection{8.1 Identity}


\begin{center}
\small
\begin{longtable}{@{}lp{10cm}@{}}
\toprule
Item & Value \\
\midrule
\endfirsthead
\toprule
Item & Value \\
\midrule
\endhead
Entry point & \texttt{spiral\_demo = boat\_mission.spiral\_demo:main} \\
Class & \texttt{SpiralDemo} \\
File & \texttt{src/boat\_mission/boat\_mission/spiral\_demo.py} \\
Launch & \texttt{spiral\_demo.launch.py} (not full mission) \\
\bottomrule
\end{longtable}
\end{center}


\subsection{8.2 Parameters}


\begin{center}
\small
\begin{longtable}{@{}llp{7.2cm}@{}}
\toprule
Name & Code default & Launch override \\
\midrule
\endfirsthead
\toprule
Name & Code default & Launch override \\
\midrule
\endhead
\texttt{pose\_topic} & \texttt{pose2d} & --- \\
\texttt{plan\_topic} & \texttt{plan} & --- \\
\texttt{center\_x} / \texttt{center\_y} & \texttt{30.0} / \texttt{20.0} & launch args (defaults −30 / −40) \\
\texttt{arrival\_radius\_m} & \texttt{8.0} & --- \\
\texttt{spiral\_ring\_spacing\_m} & \texttt{15.0} & launch \\
\texttt{spiral\_max\_radius\_m} & \texttt{80.0} & launch (e.g. 45) \\
\texttt{spiral\_step\_spacing\_m} & \texttt{10.0} & launch \\
\texttt{min\_x}…\texttt{max\_y} & ±120 & --- \\
\bottomrule
\end{longtable}
\end{center}


\textbf{No} entry in \texttt{mission.yaml} for this node.


\subsection{8.3 Subscriptions / publishers / timers}


\begin{center}
\small
\begin{longtable}{@{}lp{10cm}@{}}
\toprule
 &  \\
\midrule
\endfirsthead
\toprule
 &  \\
\midrule
\endhead
Sub & relative \texttt{pose2d} (\texttt{Pose2D}) --- track arrival at center \\
Pub & relative \texttt{plan} (\texttt{Path}), \textbf{TRANSIENT\_LOCAL} --- transit then production spiral \\
Timer & \textbf{2 Hz} (\texttt{0.5} s) phase machine \\
\bottomrule
\end{longtable}
\end{center}


\subsection{8.4 Algorithm pipeline}


\begin{itemize}[leftmargin=*]
\item Wait for pose → publish transit path to \texttt{(center\_x, center\_y)}.
\item When within \texttt{arrival\_radius\_m}, publish same expanding spiral as mission manager.
\item Stay in SPIRAL phase (no further replans).
\end{itemize}


\subsection{8.5 Dependencies}


\begin{itemize}[leftmargin=*]
\item \texttt{gazebo\_pose2d}, \texttt{los\_path\_follower} with \texttt{use\_builtin\_path:=false}; Gazebo + thruster chain. Sensing/mapping/mission usually off in demo launch.
\end{itemize}




\section{9. \texttt{boat\_navigation} --- \texttt{gazebo\_pose2d}}


\subsection{9.1 Identity}


\begin{center}
\small
\begin{longtable}{@{}lp{10cm}@{}}
\toprule
Item & Value \\
\midrule
\endfirsthead
\toprule
Item & Value \\
\midrule
\endhead
Entry point & \texttt{gazebo\_pose2d = boat\_navigation.gazebo\_pose2d:main} \\
Class & \texttt{GazeboPose2D} \\
File & \texttt{src/boat\_navigation/boat\_navigation/gazebo\_pose2d.py} \\
\bottomrule
\end{longtable}
\end{center}


\subsection{9.2 Parameters}


\begin{center}
\small
\begin{longtable}{@{}llp{7.2cm}@{}}
\toprule
Name & Code default & Launch \\
\midrule
\endfirsthead
\toprule
Name & Code default & Launch \\
\midrule
\endhead
\texttt{model\_name} & \texttt{'simple\_boat'} & \texttt{simple\_boat} / \texttt{simple\_boat\_2} \\
\texttt{world\_name} & \texttt{'niot\_world'} & \texttt{'niot\_world'} \\
\texttt{pose\_topic} & \texttt{'pose2d'} & \texttt{'pose2d'} \\
\bottomrule
\end{longtable}
\end{center}


\textbf{No} navigation.yaml entry for this node.


\subsection{9.3 I/O}


\begin{center}
\small
\begin{longtable}{@{}llp{7.2cm}@{}}
\toprule
Kind & Topic & Notes \\
\midrule
\endfirsthead
\toprule
Kind & Topic & Notes \\
\midrule
\endhead
GZ sub & \texttt{/world/\{world\_name\}/dynamic\_pose/info} & Gazebo Transport \texttt{Pose\_V} (not ROS) \\
ROS pub & relative \texttt{pose2d} & \texttt{geometry\_msgs/Pose2D} on each matching model pose \\
\bottomrule
\end{longtable}
\end{center}


\subsection{9.4 Timers / QoS}


\begin{itemize}[leftmargin=*]
\item Event-driven from Gazebo callback; ROS pub depth 10 default.
\end{itemize}


\subsection{9.5 Algorithm pipeline}


\begin{itemize}[leftmargin=*]
\item Subscribe Gazebo dynamic poses; find pose named \texttt{model\_name}.
\item Extract \(x,y\); yaw from quaternion; publish \texttt{Pose2D}.
\end{itemize}


\subsection{9.6 Dependencies}


\begin{itemize}[leftmargin=*]
\item Gazebo world \texttt{niot\_world} running with the model present. \textbf{Does not} use \texttt{ros\_gz\_bridge} for pose (direct gz transport).
\end{itemize}




\section{10. \texttt{boat\_navigation} --- \texttt{los\_path\_follower}}


\subsection{10.1 Identity}


\begin{center}
\small
\begin{longtable}{@{}lp{10cm}@{}}
\toprule
Item & Value \\
\midrule
\endfirsthead
\toprule
Item & Value \\
\midrule
\endhead
Entry point & \texttt{los\_path\_follower = boat\_navigation.los\_path\_follower:main} \\
Class & \texttt{LosPathFollower} \\
File & \texttt{src/boat\_navigation/boat\_navigation/los\_path\_follower.py} \\
Helpers & \texttt{guidance.py}, \texttt{pid.py}, \texttt{path\_utils.py} \\
\bottomrule
\end{longtable}
\end{center}


\subsection{10.2 Parameters}


All mirrored in \texttt{navigation.yaml} (\texttt{los\_path\_follower}) with same values as code defaults, except launch may set \texttt{use\_builtin\_path: false} when mission is on.


\begin{center}
\small
\begin{longtable}{@{}lp{10cm}@{}}
\toprule
Name & Default \\
\midrule
\endfirsthead
\toprule
Name & Default \\
\midrule
\endhead
\texttt{control\_rate\_hz} & \texttt{10.0} \\
\texttt{lookahead\_m} & \texttt{15.0} \\
\texttt{accept\_radius\_m} & \texttt{8.0} \\
\texttt{pass\_epsilon\_m} & \texttt{2.0} \\
\texttt{loop} & \texttt{True} \\
\texttt{u\_ref} & \texttt{0.35} \\
\texttt{u\_max} & \texttt{0.5} \\
\texttt{r\_max} & \texttt{0.3} \\
\texttt{patrol\_half\_size} & \texttt{120.0} \\
\texttt{pose\_timeout\_s} & \texttt{1.0} \\
\texttt{odom\_topic} & \texttt{odom} \\
\texttt{plan\_topic} & \texttt{plan} \\
\texttt{active\_plan\_topic} & \texttt{plan/active} \\
\texttt{cmd\_vel\_topic} & \texttt{cmd\_vel} \\
\texttt{debug\_topic} & \texttt{nav/debug} \\
\texttt{use\_builtin\_path} & \texttt{True} (false under mission) \\
\texttt{kp\_yaw}, \texttt{ki\_yaw}, \texttt{kd\_yaw}, \texttt{integral\_limit\_yaw} & \texttt{1.5}, \texttt{0.05}, \texttt{0.2}, \texttt{1.0} \\
\texttt{kp\_u}, \texttt{ki\_u}, \texttt{kd\_u}, \texttt{integral\_limit\_u} & \texttt{1.0}, \texttt{0.1}, \texttt{0.0}, \texttt{1.0} \\
\texttt{speed\_mode} & \texttt{'open\_loop'} (\texttt{'pid'} optional) \\
\bottomrule
\end{longtable}
\end{center}


\subsection{10.3 Subscriptions}


\begin{center}
\small
\begin{longtable}{@{}lllp{5.2cm}@{}}
\toprule
Topic & Rel/Abs & Type & Callback purpose \\
\midrule
\endfirsthead
\toprule
Topic & Rel/Abs & Type & Callback purpose \\
\midrule
\endhead
\texttt{odom} & relative → e.g. \texttt{/asv1/odom} & \texttt{nav\_msgs/Odometry} & Pose, yaw, measured surge; stamp for staleness. \\
\texttt{plan} & relative & \texttt{nav\_msgs/Path} & Replace active path; reindex nearest segment; \textbf{TRANSIENT\_LOCAL} sub. \\
\bottomrule
\end{longtable}
\end{center}


\subsection{10.4 Publishers}


\begin{center}
\small
\begin{longtable}{@{}lllp{5.2cm}@{}}
\toprule
Topic & Rel/Abs & Type & When \\
\midrule
\endfirsthead
\toprule
Topic & Rel/Abs & Type & When \\
\midrule
\endhead
\texttt{cmd\_vel} & relative & \texttt{Twist} & Control timer (or stop) \\
\texttt{plan/active} & relative & \texttt{Path} & When path set; \textbf{TRANSIENT\_LOCAL} \\
\texttt{nav/debug} & relative & \texttt{Float32MultiArray} & Each control tick: cross, along, heading, heading\_err, surge, segment\_index \\
\bottomrule
\end{longtable}
\end{center}


\subsection{10.5 Timers / QoS}


\begin{itemize}[leftmargin=*]
\item \textbf{Timer:} \textbf{10 Hz} control loop.
\item \textbf{QoS:} plan in/out Reliable + \textbf{TRANSIENT\_LOCAL}, depth 1.
\end{itemize}


\subsection{10.6 Algorithm pipeline}


\begin{itemize}[leftmargin=*]
\item Optional builtin square patrol if \texttt{use\_builtin\_path}.
\item If no path / stale odom → publish zero Twist.
\item LOS heading to lookahead point on current segment; PID yaw → \texttt{angular.z}.
\item Speed: open-loop scaled by heading error, or closed-loop PID on surge.
\item Advance segment on accept radius / pass geometry; loop or stop at end.
\end{itemize}


\subsection{10.7 Dependencies}


\begin{itemize}[leftmargin=*]
\item Bridge providing \texttt{odom}; \texttt{thrust\_mixer} on \texttt{cmd\_vel}; plan source = \texttt{mission\_manager} / \texttt{spiral\_demo} / builtin path.
\item Requires \texttt{autonomy:=true} in sim launch.
\end{itemize}




\section{11. \texttt{boat\_navigation} --- \texttt{trajectory\_plotter}}


\subsection{11.1 Identity}


\begin{center}
\small
\begin{longtable}{@{}lp{10cm}@{}}
\toprule
Item & Value \\
\midrule
\endfirsthead
\toprule
Item & Value \\
\midrule
\endhead
Entry point & \texttt{trajectory\_plotter = boat\_navigation.trajectory\_plotter:main} \\
Class & \texttt{TrajectoryPlotter} \\
File & \texttt{src/boat\_navigation/boat\_navigation/trajectory\_plotter.py} \\
\bottomrule
\end{longtable}
\end{center}


\subsection{11.2 Parameters}


No dedicated YAML file; \texttt{sim.launch.py} / multi-ASV block set key overlays. Code defaults:


\begin{center}
\small
\begin{longtable}{@{}llp{7.2cm}@{}}
\toprule
Name & Code default & Typical launch \\
\midrule
\endfirsthead
\toprule
Name & Code default & Typical launch \\
\midrule
\endhead
\texttt{pose\_topic} & \texttt{pose2d} & same \\
\texttt{trajectory\_topic} & \texttt{trajectory} & same \\
\texttt{active\_plan\_topic} & \texttt{plan/active} & same \\
\texttt{anomaly\_topic} & \texttt{mag/anomaly} & same \\
\texttt{asv\_namespaces} & \texttt{['']} (single relative mode) & \texttt{['asv1','asv2']} multi \\
\texttt{peak\_topic} & \texttt{/swarm/belief/peak} & same \\
\texttt{peak\_probability\_topic} & \texttt{/swarm/belief/peak\_probability} & same \\
\texttt{centroid\_topic} & \texttt{/swarm/belief/centroid} & same \\
\texttt{centroid\_spread\_topic} & \texttt{/swarm/belief/centroid\_spread} & same \\
\texttt{fix\_topic} / \texttt{fix\_rms\_topic} & \texttt{/swarm/belief/fix} / \texttt{...\_rms} & same \\
\texttt{lake\_half\_size\_m} & \texttt{150.0} & same \\
\texttt{max\_history\_points} & \texttt{20000} & --- \\
\texttt{sample\_distance\_m} & \texttt{0.2} & --- \\
\texttt{heading\_arrow\_m} & \texttt{12.0} & --- \\
\texttt{show\_true\_target} & \texttt{True} & \texttt{True} \\
\texttt{target\_x} / \texttt{target\_y} & \texttt{80.0} / \texttt{-40.0} & \textbf{\texttt{55.0} / \texttt{-35.0}} (match planted dipole) \\
\texttt{mag\_history\_points} & \texttt{6000} & --- \\
\texttt{anomaly\_vmax\_nt} & \texttt{55.0} & \texttt{55.0} \\
\texttt{anomaly\_hit\_threshold\_nt} & \texttt{15.0} & \texttt{15.0} \\
\texttt{mission\_state\_topic} & \texttt{mission/state} & --- \\
\texttt{verify\_result\_topic} & \texttt{/swarm/verify/result} & --- \\
\texttt{mission\_complete\_topic} & \texttt{/swarm/mission/complete} & --- \\
\texttt{mission\_status\_topic} & \texttt{/swarm/mission/status} & --- \\
\texttt{verify\_confirmations\_required} & \texttt{4} & --- \\
\bottomrule
\end{longtable}
\end{center}


\subsection{11.3 Subscriptions (per boat + swarm)}


\textbf{Per boat} (relative, or \texttt{/ns/...} when \texttt{asv\_namespaces} set):


\begin{center}
\small
\begin{longtable}{@{}llp{7.2cm}@{}}
\toprule
Topic & Type & Purpose \\
\midrule
\endfirsthead
\toprule
Topic & Type & Purpose \\
\midrule
\endhead
\texttt{pose2d} & \texttt{Pose2D} & Track history / heading \\
\texttt{plan/active} & \texttt{Path} & Draw LOS route (\textbf{TRANSIENT\_LOCAL}) \\
\texttt{mag/anomaly} & \texttt{MagAnomaly} & Mag scatter + time series \\
\texttt{mission/state} & \texttt{MissionState} & Phase banner \\
\bottomrule
\end{longtable}
\end{center}


\textbf{Swarm / absolute:}


\begin{center}
\small
\begin{longtable}{@{}llp{7.2cm}@{}}
\toprule
Topic & Type & Purpose \\
\midrule
\endfirsthead
\toprule
Topic & Type & Purpose \\
\midrule
\endhead
peak / peak\_probability & \texttt{PoseStamped} / \texttt{Float64} & Peak overlay \\
centroid / centroid\_spread & \texttt{PoseStamped} / \texttt{Float64} & Estimate + σ ring \\
fix / fix\_rms & \texttt{PoseStamped} / \texttt{Float64} & Dipole fix \\
\texttt{/swarm/verify/result} & \texttt{VerifyResult} & Confirmation UI \\
\texttt{/swarm/mission/complete} & \texttt{Bool} & Latch confirmed (\textbf{TL}) \\
\texttt{/swarm/mission/status} & \texttt{String} & Status text (\textbf{TL}) \\
\bottomrule
\end{longtable}
\end{center}


\subsection{11.4 Publishers / timers}


\begin{center}
\small
\begin{longtable}{@{}lp{10cm}@{}}
\toprule
 &  \\
\midrule
\endfirsthead
\toprule
 &  \\
\midrule
\endhead
Pub & \texttt{\{ns\}/trajectory} or relative \texttt{trajectory} (\texttt{Path}), \textbf{TRANSIENT\_LOCAL}, \textbf{1 Hz} \\
Timer & \textbf{5 Hz} (\texttt{0.2} s) matplotlib redraw; \textbf{1 Hz} trajectory publish \\
\bottomrule
\end{longtable}
\end{center}


\subsection{11.5 Algorithm pipeline}


\begin{itemize}[leftmargin=*]
\item Maintain per-boat track, route, anomaly history.
\item Overlay true target, centroid, peak cell, dipole fix, shore box.
\item Status strip maps modes → SEARCHING / HUNTING / HOLD / VERIFYING / CONFIRMED.
\item Visualization only --- no control outputs.
\end{itemize}


\subsection{11.6 Dependencies}


\begin{itemize}[leftmargin=*]
\item Optional: any combination of pose, plan, anomaly, belief, mission topics. Disabled when \texttt{headless} or \texttt{plot\_trajectory:=false}.
\end{itemize}




\chapter{Part IV --- Executable Index}
\label{ch:ros-catalog}

This chapter is the authoritative software reference for every message, bridge, and node in the dual-ASV stack. Each node section lists identity, parameters, subscriptions, publications, timers, QoS, internal algorithm steps, and runtime dependencies.


\begin{center}
\small
\begin{longtable}{@{}lllp{5.2cm}@{}}
\toprule
Package & Console script & Class & Source file \\
\midrule
\endfirsthead
\toprule
Package & Console script & Class & Source file \\
\midrule
\endhead
\texttt{boat\_control} & \texttt{thrust\_mixer} & \texttt{ThrustMixer} & \texttt{boat\_control/mixer.py} \\
\texttt{boat\_sensing} & \texttt{mag\_driver} & \texttt{MagDriver} & \texttt{boat\_sensing/mag\_driver.py} \\
\texttt{boat\_sensing} & \texttt{mag\_filter} & \texttt{MagFilter} & \texttt{boat\_sensing/mag\_filter.py} \\
\texttt{boat\_sensing} & \texttt{calibration\_node} & \texttt{CalibrationNode} & \texttt{boat\_sensing/calibration\_node.py} \\
\texttt{boat\_mapping} & \texttt{bayes\_fusion} & \texttt{BayesFusionNode} & \texttt{boat\_mapping/bayes\_fusion.py} \\
\texttt{boat\_mission} & \texttt{mission\_manager} & \texttt{MissionManager} & \texttt{boat\_mission/mission\_manager.py} \\
\texttt{boat\_mission} & \texttt{verify\_coordinator} & \texttt{VerifyCoordinator} & \texttt{boat\_mission/verify\_coordinator.py} \\
\texttt{boat\_mission} & \texttt{spiral\_demo} & \texttt{SpiralDemo} & \texttt{boat\_mission/spiral\_demo.py} \\
\texttt{boat\_navigation} & \texttt{gazebo\_pose2d} & \texttt{GazeboPose2D} & \texttt{boat\_navigation/gazebo\_pose2d.py} \\
\texttt{boat\_navigation} & \texttt{los\_path\_follower} & \texttt{LosPathFollower} & \texttt{boat\_navigation/los\_path\_follower.py} \\
\texttt{boat\_navigation} & \texttt{trajectory\_plotter} & \texttt{TrajectoryPlotter} & \texttt{boat\_navigation/trajectory\_plotter.py} \\
\texttt{ros\_gz\_bridge} & \texttt{parameter\_bridge} (\texttt{boat\_bridge}) & (C++) & \texttt{config/bridge.yaml} \\
\texttt{ros\_gz\_bridge} & \texttt{parameter\_bridge} (\texttt{clock\_bridge}) & (C++) & \texttt{config/clock\_bridge.yaml} \\
\bottomrule
\end{longtable}
\end{center}




\chapter{Part V --- QoS Summary (non-default)}
\label{ch:ros-catalog}

This chapter is the authoritative software reference for every message, bridge, and node in the dual-ASV stack. Each node section lists identity, parameters, subscriptions, publications, timers, QoS, internal algorithm steps, and runtime dependencies.


\begin{center}
\small
\begin{longtable}{@{}llp{7.2cm}@{}}
\toprule
Node & Topic(s) & Policy \\
\midrule
\endfirsthead
\toprule
Node & Topic(s) & Policy \\
\midrule
\endhead
\texttt{mission\_manager} & \texttt{plan}, \texttt{/swarm/verify/result} & Reliable + TRANSIENT\_LOCAL, depth 1 \\
\texttt{verify\_coordinator} & \texttt{/swarm/mission/complete}, \texttt{halt}, \texttt{/swarm/mission/status} & Reliable + TRANSIENT\_LOCAL, depth 1 \\
\texttt{spiral\_demo} & \texttt{plan} & Reliable + TRANSIENT\_LOCAL, depth 1 \\
\texttt{los\_path\_follower} & \texttt{plan} (sub), \texttt{plan/active} (pub) & Reliable + TRANSIENT\_LOCAL, depth 1 \\
\texttt{trajectory\_plotter} & \texttt{plan/active} (sub), \texttt{trajectory} (pub), complete/status (sub) & Reliable + TRANSIENT\_LOCAL, depth 1 \\
\bottomrule
\end{longtable}
\end{center}


All other pubs/subs use rclpy defaults (typically Reliable, Volatile, depth as coded 10/20/50).




\chapter{Part VI --- Minimum Runtime Graph (single ASV, full stack)}
\label{ch:ros-catalog}

This chapter is the authoritative software reference for every message, bridge, and node in the dual-ASV stack. Each node section lists identity, parameters, subscriptions, publications, timers, QoS, internal algorithm steps, and runtime dependencies.


Must be running for a closed-loop magnetic search demo:


\begin{enumerate}[leftmargin=*]
\item Gazebo (\texttt{niot\_world} + \texttt{simple\_boat})
\item \texttt{boat\_bridge} (+ \texttt{clock\_bridge} if \texttt{use\_sim\_time})
\item \texttt{/asv1/thrust\_mixer}
\item \texttt{/asv1/gazebo\_pose2d}
\item \texttt{/asv1/los\_path\_follower} (\texttt{use\_builtin\_path:=false})
\item \texttt{/asv1/mag\_driver} → \texttt{mag\_filter} → \texttt{calibration\_node}
\item \texttt{bayes\_fusion}
\item \texttt{/asv1/mission\_manager}
\item \texttt{verify\_coordinator} (self-verify or multi-ASV)
\item Optional: \texttt{trajectory\_plotter}
\end{enumerate}


Two-ASV adds mirrored \texttt{/asv2/*} stack, dual anomaly topics into fusion, \texttt{self\_verify:=false}, and peer verification via coordinator.




\emph{End of catalog. Generated from source under \texttt{version\_3/src}.}

\chapter{Launches, configuration, and operations}
\label{ch:ops}

\section{Primary launches}
\begin{center}
\begin{tabular}{@{}lp{8.8cm}@{}}
\toprule
Launch & Purpose \\
\midrule
\texttt{sim.launch.py} & Core world + \(1\)--\(2\) ASV stacks, feature flags \\
\texttt{multi\_asv\_phase1.launch.py} & Two boats, no autonomy/sensing \\
\texttt{multi\_asv\_phase2.launch.py} & Voronoi lawnmower only \\
\texttt{multi\_asv\_phase3.launch.py} & Full stack + shared belief \\
\texttt{multi\_asv\_phase4.launch.py} & Cooperative verify (demo default) \\
\texttt{record.launch.py} / \texttt{replay.launch.py} & Bag workflows \\
\bottomrule
\end{tabular}
\end{center}

\section{YAML files (\texttt{boat\_bringup/config})}
\begin{itemize}[leftmargin=*]
  \item \texttt{sensing.yaml} --- driver, filter, calibration, planted dipole.
  \item \texttt{mapping.yaml} --- belief grid, HIT thresholds, dipole fit.
  \item \texttt{mission.yaml} --- coverage box, lawnmower, FSM gates, verify.
  \item \texttt{navigation.yaml} --- LOS / PID / surge.
  \item \texttt{bridge.yaml}, \texttt{clock\_bridge.yaml} --- Gazebo bridges.
\end{itemize}
Launch remaps and per-ASV overrides (namespace, model name, lawnmower index,
\texttt{self\_verify}) sit in \texttt{sim.launch.py}.

\section{Fast mode}
\texttt{fast:=true} raises Gazebo RTF, starts \texttt{clock\_bridge}, and sets
\texttt{use\_sim\_time} so timers follow \texttt{/clock}.

\section{Inspecting the live graph}
\begin{verbatim}
ros2 node list
ros2 topic list | grep -E 'asv|swarm'
ros2 topic echo /asv1/mission/state --once
ros2 topic echo /swarm/belief/fix --once
ros2 topic echo /swarm/mission/status --once
\end{verbatim}

\section{Build and run (\texttt{version\_3})}
\begin{verbatim}
cd version_3
source /opt/ros/jazzy/setup.bash
colcon build --symlink-install
source install/setup.bash
ros2 launch boat_bringup multi_asv_phase4.launch.py
\end{verbatim}

\chapter{Acceptance, oscillation, and scoring}
\label{ch:accept}

\section{There is no onboard ``error $<5\,\mathrm{m}$ countdown''}
Robots do not receive the planted ground-truth pose. Completion is gated by:
\begin{enumerate}[leftmargin=*]
  \item consecutive high peak probability (enter hunt);
  \item spiral minimum duration (declare);
  \item verification confirmation count (COMPLETE);
  \item optional dipole residual quality (publish fix).
\end{enumerate}

\section{Why the estimate oscillates then settles}
The discrete peak lives on a \(10\,\mathrm{m}\) grid. Early HITs create a ridge
of plausible cells; \(\arg\max\) jumps among neighbours. The centroid is
smoother; the dipole least-squares fix usually becomes the stable few-metre
estimate after spiral/verify angular diversity.

\section{How to score $<5\,\mathrm{m}$ for a competition}
Offline (or judge script), compare
\texttt{/swarm/belief/fix} (and centroid) to the planted
\((t_x,t_y)\). Report final error at COMPLETE and time spent below
\(5\,\mathrm{m}\). Do not claim the vehicle waited on that threshold unless you
add an external scoring node.

\section{Operator confirmation UX}
When verify succeeds, the plotter status strip switches to green
\textbf{CONFIRMED} and shows verifier id and confirmation count --- that is the
human-visible ``yes, verification phase finished successfully'' signal.


\backmatter
\begin{thebibliography}{99}
\bibitem{ros2} Open Robotics, \emph{ROS~2 Jazzy Documentation},
  \url{https://docs.ros.org/en/jazzy/}.
\bibitem{gazebo} Open Robotics, \emph{Gazebo Sim Harmonic},
  \url{https://gazebosim.org/}.
\bibitem{shannon} C.~E.~Shannon, ``A mathematical theory of communication,''
  \emph{Bell Syst.\ Tech.\ J.}, 1948.
\bibitem{lm} K.~Levenberg, ``A method for the solution of certain non-linear
  problems in least squares,'' \emph{Q.\ Appl.\ Math.}, 1944.
\end{thebibliography}

\end{document}
